% The Moebius loop: a delay fed back with a sign flip, three times over: the circuit you build, the loop it is, the sound you hear.
% The heavy stroke is the path that comes back: the feedback wire in the circuit, the path round the band on its own.
% Data and geometry: mobius-loop-odd-half-series.py (mlohs-comb.dat, mlohs-peaks.dat, mobius-strip.tex). Everything here is computed (ink).
% Compile: lualatex mobius-loop-odd-half-series.tex ; SVG via PyMuPDF (dvisvgm's PDF route is broken on this Mac).
\documentclass[border=0pt]{standalone}
\usepackage{fontspec}
\setmainfont{JetBrains Mono}
\setmonofont{JetBrains Mono}
\usepackage{tikz}
\usepackage{pgfplots}
\usepackage{array}
\pgfplotsset{compat=1.16}
\usetikzlibrary{arrows.meta, calc, backgrounds, decorations.markings, positioning}
% Drafting Table, paper ground (_ops/drafting-table/design-system/colors_and_type.css)
\definecolor{paper}{HTML}{F0EEE5}
\definecolor{ink}{HTML}{1A1916}
\definecolor{ink2}{HTML}{3A3A35}
\definecolor{ink3}{HTML}{6A6862}
\definecolor{slab}{HTML}{D8D4C6}
\tikzset{
  >={Stealth[length=1.8mm]},
  loop/.style={draw=ink, line width=1.5pt},
  back/.style={draw=ink, line width=1.5pt, ->},
  wire/.style={draw=ink, line width=0.6pt},
  blk/.style={rectangle, draw=ink, fill=slab, line width=0.7pt, minimum height=7mm, minimum width=12mm, inner sep=2pt, font=\scriptsize, text=ink},
  adder/.style={circle, draw=ink, fill=paper, line width=0.7pt, minimum size=4.6mm, inner sep=0pt, font=\scriptsize},
  tap/.style={circle, fill=ink, minimum size=3pt, inner sep=0pt},
  lab/.style={font=\scriptsize, text=ink2},
  rowhead/.style={font=\bfseries\small, anchor=west, text=ink},
  rowsub/.style={font=\scriptsize, anchor=west, text=ink2},
  colhead/.style={font=\bfseries\scriptsize, anchor=west, text=ink},
  spin/.style={postaction={decorate}, decoration={markings, mark=at position 0.25 with {\arrow{>}}}},
}
\pgfplotsset{
  /pgf/number format/assume math mode=true,
  snd/.style={width=6.6cm, height=3.0cm, axis lines=left, tick align=outside, ymin=0, ytick=\empty,
    axis line style={ink, line width=0.5pt}, tick style={ink}, clip=false,
    xlabel style={font=\scriptsize, text=ink2, yshift=1mm}, ylabel style={font=\scriptsize, text=ink2},
    tick label style={font=\tiny, text=ink2}},
}
% a row: frame (top y, height), heading, subtitle, and the two column rules
\newcommand{\rowframe}[4]{%
  \draw[ink, line width=0.5pt] (-0.3,#1) rectangle (19.9,#1-#2);
  \draw[ink, line width=0.3pt] (6.9,#1-0.95) -- (6.9,#1-#2) (12.4,#1-0.95) -- (12.4,#1-#2) (-0.3,#1-0.95) -- (19.9,#1-0.95);
  \node[rowhead] at (0,#1-0.32) {#3};
  \node[rowsub] at (0,#1-0.7) {#4};}
\begin{document}
\begin{tikzpicture}

% ---------- column heads ----------
\node[colhead] at (0,1.15) {THE CIRCUIT \quad what you build};
\node[colhead] at (7.2,1.15) {THE LOOP \quad what it is};
\node[colhead] at (12.7,1.15) {THE SOUND \quad what you hear};
\node[font=\scriptsize, anchor=west, text=ink2] at (0,0.55) {Heavy: the path that comes back. In the circuit it is the feedback wire; on its own it is a path round a band with one half twist.};

% =============== THE MÖBIUS LOOP ===============
\def\y{0}
\rowframe{\y}{6.6}{THE MÖBIUS LOOP}{flip the sign each trip: the odd multiples of half the loop's frequency, 120 · 360 · 600 · 840 Hz}
\begin{scope}[shift={(0,\y)}]
  % the circuit
  \node[lab] (in) at (0.1,-3.0) {in};
  \node[adder] (a) at (0.9,-3.0) {+};
  \node[blk] (d) at (2.6,-3.0) {delay T};
  \node[tap] (t) at (4.3,-3.0) {};
  \draw[wire,->] (in) -- (a); \draw[wire,->] (a) -- (d); \draw[wire] (d) -- (t);
  \draw[wire,->] (t) -- (6.65,-3.0) node[lab, above left] {out};
  \node[blk, minimum width=12mm, minimum height=8mm] (g) at (2.6,-4.6) {× {\large\bfseries −a}};
  \draw[back] (t) |- (g.east);
  \draw[back] (g.west) -| (a);
  \node[lab, anchor=north, align=center] at (2.75,-5.25) {what comes back is turned\\upside down every trip};

  % the loop: a band with one half twist, and a plain circle for contrast
  \begin{scope}[shift={(8.7,-3.05)}, scale=1.15]
    % a Moebius band with a path that needs two laps to close; generated by mobius-loop-odd-half-series.py
\filldraw[fill=slab!27!paper, draw=slab!27!paper, line width=0.2pt] (0.035,0.170) -- (0.000,0.174) -- (0.000,0.886) -- (0.035,0.890) -- cycle;
\filldraw[fill=slab!27!paper, draw=slab!27!paper, line width=0.2pt] (0.000,0.174) -- (-0.035,0.177) -- (-0.035,0.882) -- (0.000,0.886) -- cycle;
\draw[ink2, line width=0.6pt, line cap=round] (0.035,0.170) -- (0.000,0.174);
\draw[ink2, line width=0.6pt, line cap=round] (0.035,0.890) -- (0.000,0.886);
\draw[ink2, line width=0.6pt, line cap=round] (0.000,0.174) -- (-0.035,0.177);
\draw[ink2, line width=0.6pt, line cap=round] (0.000,0.886) -- (-0.035,0.882);
\draw[ink3!50, line width=0.25pt] (0.000,0.174) -- (0.000,0.886);
\draw[loop, line cap=round] (0.035,0.710) -- (0.000,0.708);
\draw[loop, line cap=round] (0.000,0.708) -- (-0.035,0.706);
\filldraw[fill=slab!26!paper, draw=slab!26!paper, line width=0.2pt] (0.069,0.165) -- (0.035,0.170) -- (0.035,0.890) -- (0.071,0.892) -- cycle;
\draw[ink2, line width=0.6pt, line cap=round] (0.069,0.165) -- (0.035,0.170);
\draw[ink2, line width=0.6pt, line cap=round] (0.071,0.892) -- (0.035,0.890);
\draw[loop, line cap=round] (0.070,0.710) -- (0.035,0.710);
\filldraw[fill=slab!28!paper, draw=slab!28!paper, line width=0.2pt] (-0.035,0.177) -- (-0.071,0.180) -- (-0.069,0.877) -- (-0.035,0.882) -- cycle;
\draw[ink2, line width=0.6pt, line cap=round] (-0.035,0.177) -- (-0.071,0.180);
\draw[ink2, line width=0.6pt, line cap=round] (-0.035,0.882) -- (-0.069,0.877);
\draw[loop, line cap=round] (-0.035,0.706) -- (-0.069,0.703);
\draw[loop, line cap=round] (-0.035,0.353) -- (-0.070,0.355);
\filldraw[fill=slab!26!paper, draw=slab!26!paper, line width=0.2pt] (0.102,0.160) -- (0.069,0.165) -- (0.071,0.892) -- (0.107,0.894) -- cycle;
\filldraw[fill=slab!28!paper, draw=slab!28!paper, line width=0.2pt] (-0.071,0.180) -- (-0.107,0.183) -- (-0.102,0.871) -- (-0.069,0.877) -- cycle;
\draw[ink2, line width=0.6pt, line cap=round] (0.102,0.160) -- (0.069,0.165);
\draw[ink2, line width=0.6pt, line cap=round] (0.107,0.894) -- (0.071,0.892);
\draw[ink2, line width=0.6pt, line cap=round] (-0.071,0.180) -- (-0.107,0.183);
\draw[ink2, line width=0.6pt, line cap=round] (-0.069,0.877) -- (-0.102,0.871);
\draw[loop, line cap=round] (0.106,0.711) -- (0.070,0.710);
\draw[loop, line cap=round] (0.103,0.343) -- (0.069,0.347);
\draw[loop, line cap=round] (-0.069,0.703) -- (-0.103,0.699);
\draw[loop, line cap=round] (-0.070,0.355) -- (-0.106,0.355);
\filldraw[fill=slab!26!paper, draw=slab!26!paper, line width=0.2pt] (0.135,0.154) -- (0.102,0.160) -- (0.107,0.894) -- (0.143,0.895) -- cycle;
\filldraw[fill=slab!29!paper, draw=slab!29!paper, line width=0.2pt] (-0.107,0.183) -- (-0.143,0.185) -- (-0.135,0.865) -- (-0.102,0.871) -- cycle;
\draw[ink2, line width=0.6pt, line cap=round] (0.135,0.154) -- (0.102,0.160);
\draw[ink2, line width=0.6pt, line cap=round] (0.143,0.895) -- (0.107,0.894);
\draw[ink2, line width=0.6pt, line cap=round] (-0.107,0.183) -- (-0.143,0.185);
\draw[ink2, line width=0.6pt, line cap=round] (-0.102,0.871) -- (-0.135,0.865);
\draw[loop, line cap=round] (0.141,0.710) -- (0.106,0.711);
\draw[loop, line cap=round] (0.137,0.339) -- (0.103,0.343);
\draw[loop, line cap=round] (-0.103,0.699) -- (-0.137,0.695);
\draw[loop, line cap=round] (-0.106,0.355) -- (-0.141,0.355);
\filldraw[fill=slab!26!paper, draw=slab!26!paper, line width=0.2pt] (0.167,0.148) -- (0.135,0.154) -- (0.143,0.895) -- (0.180,0.896) -- cycle;
\filldraw[fill=slab!29!paper, draw=slab!29!paper, line width=0.2pt] (-0.143,0.185) -- (-0.180,0.186) -- (-0.167,0.858) -- (-0.135,0.865) -- cycle;
\draw[ink2, line width=0.6pt, line cap=round] (0.167,0.148) -- (0.135,0.154);
\draw[ink2, line width=0.6pt, line cap=round] (0.180,0.896) -- (0.143,0.895);
\draw[ink2, line width=0.6pt, line cap=round] (-0.143,0.185) -- (-0.180,0.186);
\draw[ink2, line width=0.6pt, line cap=round] (-0.135,0.865) -- (-0.167,0.858);
\draw[loop, line cap=round] (0.177,0.709) -- (0.141,0.710);
\draw[loop, line cap=round] (0.170,0.335) -- (0.137,0.339);
\draw[loop, line cap=round] (-0.137,0.695) -- (-0.170,0.690);
\draw[loop, line cap=round] (-0.141,0.355) -- (-0.177,0.354);
\filldraw[fill=slab!25!paper, draw=slab!25!paper, line width=0.2pt] (0.199,0.141) -- (0.167,0.148) -- (0.180,0.896) -- (0.217,0.895) -- cycle;
\draw[ink2, line width=0.6pt, line cap=round] (0.199,0.141) -- (0.167,0.148);
\draw[ink2, line width=0.6pt, line cap=round] (0.217,0.895) -- (0.180,0.896);
\draw[loop, line cap=round] (0.212,0.707) -- (0.177,0.709);
\draw[loop, line cap=round] (0.203,0.330) -- (0.170,0.335);
\filldraw[fill=slab!30!paper, draw=slab!30!paper, line width=0.2pt] (-0.180,0.186) -- (-0.217,0.187) -- (-0.199,0.850) -- (-0.167,0.858) -- cycle;
\draw[ink2, line width=0.6pt, line cap=round] (-0.180,0.186) -- (-0.217,0.187);
\draw[ink2, line width=0.6pt, line cap=round] (-0.167,0.858) -- (-0.199,0.850);
\draw[loop, line cap=round] (-0.170,0.690) -- (-0.203,0.684);
\filldraw[fill=slab!25!paper, draw=slab!25!paper, line width=0.2pt] (0.230,0.134) -- (0.199,0.141) -- (0.217,0.895) -- (0.254,0.894) -- cycle;
\filldraw[fill=slab!31!paper, draw=slab!31!paper, line width=0.2pt] (-0.217,0.187) -- (-0.254,0.187) -- (-0.230,0.841) -- (-0.199,0.850) -- cycle;
\draw[ink2, line width=0.6pt, line cap=round] (0.230,0.134) -- (0.199,0.141);
\draw[ink2, line width=0.6pt, line cap=round] (0.254,0.894) -- (0.217,0.895);
\draw[ink2, line width=0.6pt, line cap=round] (-0.217,0.187) -- (-0.254,0.187);
\draw[ink2, line width=0.6pt, line cap=round] (-0.199,0.850) -- (-0.230,0.841);
\draw[loop, line cap=round] (0.248,0.704) -- (0.212,0.707);
\draw[loop, line cap=round] (-0.203,0.684) -- (-0.236,0.678);
\filldraw[fill=slab!25!paper, draw=slab!25!paper, line width=0.2pt] (0.260,0.127) -- (0.230,0.134) -- (0.254,0.894) -- (0.292,0.892) -- cycle;
\filldraw[fill=slab!32!paper, draw=slab!32!paper, line width=0.2pt] (-0.254,0.187) -- (-0.292,0.186) -- (-0.260,0.832) -- (-0.230,0.841) -- cycle;
\draw[ink2, line width=0.6pt, line cap=round] (0.260,0.127) -- (0.230,0.134);
\draw[ink2, line width=0.6pt, line cap=round] (0.292,0.892) -- (0.254,0.894);
\draw[ink2, line width=0.6pt, line cap=round] (-0.254,0.187) -- (-0.292,0.186);
\draw[ink2, line width=0.6pt, line cap=round] (-0.230,0.841) -- (-0.260,0.832);
\draw[loop, line cap=round] (0.284,0.701) -- (0.248,0.704);
\draw[loop, line cap=round] (-0.236,0.678) -- (-0.268,0.671);
\filldraw[fill=slab!32!paper, draw=slab!32!paper, line width=0.2pt] (-0.292,0.186) -- (-0.329,0.185) -- (-0.289,0.823) -- (-0.260,0.832) -- cycle;
\draw[ink2, line width=0.6pt, line cap=round] (-0.292,0.186) -- (-0.329,0.185);
\draw[ink2, line width=0.6pt, line cap=round] (-0.260,0.832) -- (-0.289,0.823);
\draw[loop, line cap=round] (-0.268,0.671) -- (-0.299,0.663);
\draw[loop, line cap=round] (-0.284,0.348) -- (-0.319,0.345);
\filldraw[fill=slab!25!paper, draw=slab!25!paper, line width=0.2pt] (0.289,0.119) -- (0.260,0.127) -- (0.292,0.892) -- (0.329,0.889) -- cycle;
\draw[ink2, line width=0.6pt, line cap=round] (0.289,0.119) -- (0.260,0.127);
\draw[ink2, line width=0.6pt, line cap=round] (0.329,0.889) -- (0.292,0.892);
\draw[loop, line cap=round] (0.319,0.696) -- (0.284,0.701);
\filldraw[fill=slab!33!paper, draw=slab!33!paper, line width=0.2pt] (-0.329,0.185) -- (-0.367,0.184) -- (-0.317,0.812) -- (-0.289,0.823) -- cycle;
\draw[ink2, line width=0.6pt, line cap=round] (-0.329,0.185) -- (-0.367,0.184);
\draw[ink2, line width=0.6pt, line cap=round] (-0.289,0.823) -- (-0.317,0.812);
\draw[loop, line cap=round] (-0.299,0.663) -- (-0.330,0.655);
\draw[loop, line cap=round] (-0.319,0.345) -- (-0.354,0.341);
\filldraw[fill=slab!25!paper, draw=slab!25!paper, line width=0.2pt] (0.317,0.111) -- (0.289,0.119) -- (0.329,0.889) -- (0.367,0.885) -- cycle;
\draw[ink2, line width=0.6pt, line cap=round] (0.317,0.111) -- (0.289,0.119);
\draw[ink2, line width=0.6pt, line cap=round] (0.367,0.885) -- (0.329,0.889);
\draw[loop, line cap=round] (0.354,0.692) -- (0.319,0.696);
\draw[loop, line cap=round] (0.330,0.304) -- (0.299,0.312);
\filldraw[fill=slab!34!paper, draw=slab!34!paper, line width=0.2pt] (-0.367,0.184) -- (-0.405,0.181) -- (-0.345,0.802) -- (-0.317,0.812) -- cycle;
\draw[ink2, line width=0.6pt, line cap=round] (-0.367,0.184) -- (-0.405,0.181);
\draw[ink2, line width=0.6pt, line cap=round] (-0.317,0.812) -- (-0.345,0.802);
\draw[loop, line cap=round] (-0.330,0.655) -- (-0.360,0.646);
\draw[loop, line cap=round] (-0.354,0.341) -- (-0.390,0.336);
\filldraw[fill=slab!26!paper, draw=slab!26!paper, line width=0.2pt] (0.345,0.102) -- (0.317,0.111) -- (0.367,0.885) -- (0.405,0.880) -- cycle;
\draw[ink2, line width=0.6pt, line cap=round] (0.345,0.102) -- (0.317,0.111);
\draw[ink2, line width=0.6pt, line cap=round] (0.405,0.880) -- (0.367,0.885);
\draw[loop, line cap=round] (0.390,0.686) -- (0.354,0.692);
\draw[loop, line cap=round] (0.360,0.297) -- (0.330,0.304);
\filldraw[fill=slab!35!paper, draw=slab!35!paper, line width=0.2pt] (-0.405,0.181) -- (-0.442,0.178) -- (-0.371,0.790) -- (-0.345,0.802) -- cycle;
\draw[ink2, line width=0.6pt, line cap=round] (-0.405,0.181) -- (-0.442,0.178);
\draw[ink2, line width=0.6pt, line cap=round] (-0.345,0.802) -- (-0.371,0.790);
\draw[loop, line cap=round] (-0.360,0.646) -- (-0.389,0.637);
\draw[loop, line cap=round] (-0.390,0.336) -- (-0.424,0.331);
\filldraw[fill=slab!26!paper, draw=slab!26!paper, line width=0.2pt] (0.371,0.093) -- (0.345,0.102) -- (0.405,0.880) -- (0.442,0.875) -- cycle;
\draw[ink2, line width=0.6pt, line cap=round] (0.371,0.093) -- (0.345,0.102);
\draw[ink2, line width=0.6pt, line cap=round] (0.442,0.875) -- (0.405,0.880);
\draw[loop, line cap=round] (0.424,0.679) -- (0.390,0.686);
\draw[loop, line cap=round] (0.389,0.289) -- (0.360,0.297);
\filldraw[fill=slab!26!paper, draw=slab!26!paper, line width=0.2pt] (0.397,0.084) -- (0.371,0.093) -- (0.442,0.875) -- (0.480,0.868) -- cycle;
\filldraw[fill=slab!36!paper, draw=slab!36!paper, line width=0.2pt] (-0.442,0.178) -- (-0.480,0.174) -- (-0.397,0.778) -- (-0.371,0.790) -- cycle;
\draw[ink2, line width=0.6pt, line cap=round] (0.397,0.084) -- (0.371,0.093);
\draw[ink2, line width=0.6pt, line cap=round] (0.480,0.868) -- (0.442,0.875);
\draw[ink2, line width=0.6pt, line cap=round] (-0.442,0.178) -- (-0.480,0.174);
\draw[ink2, line width=0.6pt, line cap=round] (-0.371,0.790) -- (-0.397,0.778);
\draw[loop, line cap=round] (0.459,0.672) -- (0.424,0.679);
\draw[loop, line cap=round] (0.418,0.280) -- (0.389,0.289);
\draw[loop, line cap=round] (-0.389,0.637) -- (-0.418,0.627);
\filldraw[fill=slab!26!paper, draw=slab!26!paper, line width=0.2pt] (0.422,0.075) -- (0.397,0.084) -- (0.480,0.868) -- (0.517,0.861) -- cycle;
\filldraw[fill=slab!38!paper, draw=slab!38!paper, line width=0.2pt] (-0.480,0.174) -- (-0.517,0.170) -- (-0.422,0.766) -- (-0.397,0.778) -- cycle;
\draw[ink2, line width=0.6pt, line cap=round] (0.422,0.075) -- (0.397,0.084);
\draw[ink2, line width=0.6pt, line cap=round] (0.517,0.861) -- (0.480,0.868);
\draw[ink2, line width=0.6pt, line cap=round] (-0.480,0.174) -- (-0.517,0.170);
\draw[ink2, line width=0.6pt, line cap=round] (-0.397,0.778) -- (-0.422,0.766);
\draw[loop, line cap=round] (0.493,0.664) -- (0.459,0.672);
\draw[loop, line cap=round] (-0.418,0.627) -- (-0.446,0.617);
\filldraw[fill=slab!27!paper, draw=slab!27!paper, line width=0.2pt] (0.446,0.065) -- (0.422,0.075) -- (0.517,0.861) -- (0.554,0.853) -- cycle;
\draw[ink2, line width=0.6pt, line cap=round] (0.446,0.065) -- (0.422,0.075);
\draw[ink2, line width=0.6pt, line cap=round] (0.554,0.853) -- (0.517,0.861);
\draw[ink3!50, line width=0.25pt] (0.446,0.065) -- (0.554,0.853);
\draw[loop, line cap=round] (0.527,0.656) -- (0.493,0.664);
\filldraw[fill=slab!39!paper, draw=slab!39!paper, line width=0.2pt] (-0.517,0.170) -- (-0.554,0.165) -- (-0.446,0.753) -- (-0.422,0.766) -- cycle;
\draw[ink2, line width=0.6pt, line cap=round] (-0.517,0.170) -- (-0.554,0.165);
\draw[ink2, line width=0.6pt, line cap=round] (-0.422,0.766) -- (-0.446,0.753);
\draw[loop, line cap=round] (-0.446,0.617) -- (-0.473,0.606);
\filldraw[fill=slab!27!paper, draw=slab!27!paper, line width=0.2pt] (0.469,0.055) -- (0.446,0.065) -- (0.554,0.853) -- (0.591,0.844) -- cycle;
\draw[ink2, line width=0.6pt, line cap=round] (0.469,0.055) -- (0.446,0.065);
\draw[ink2, line width=0.6pt, line cap=round] (0.591,0.844) -- (0.554,0.853);
\draw[loop, line cap=round] (0.561,0.647) -- (0.527,0.656);
\filldraw[fill=slab!40!paper, draw=slab!40!paper, line width=0.2pt] (-0.554,0.165) -- (-0.591,0.159) -- (-0.469,0.740) -- (-0.446,0.753) -- cycle;
\draw[ink2, line width=0.6pt, line cap=round] (-0.554,0.165) -- (-0.591,0.159);
\draw[ink2, line width=0.6pt, line cap=round] (-0.446,0.753) -- (-0.469,0.740);
\draw[ink3!50, line width=0.25pt] (-0.554,0.165) -- (-0.446,0.753);
\draw[loop, line cap=round] (-0.473,0.606) -- (-0.499,0.595);
\draw[loop, line cap=round] (-0.527,0.312) -- (-0.561,0.304);
\filldraw[fill=slab!28!paper, draw=slab!28!paper, line width=0.2pt] (0.491,0.045) -- (0.469,0.055) -- (0.591,0.844) -- (0.628,0.834) -- cycle;
\filldraw[fill=slab!41!paper, draw=slab!41!paper, line width=0.2pt] (-0.591,0.159) -- (-0.628,0.153) -- (-0.491,0.726) -- (-0.469,0.740) -- cycle;
\draw[ink2, line width=0.6pt, line cap=round] (0.491,0.045) -- (0.469,0.055);
\draw[ink2, line width=0.6pt, line cap=round] (0.628,0.834) -- (0.591,0.844);
\draw[ink2, line width=0.6pt, line cap=round] (-0.591,0.159) -- (-0.628,0.153);
\draw[ink2, line width=0.6pt, line cap=round] (-0.469,0.740) -- (-0.491,0.726);
\draw[loop, line cap=round] (0.594,0.637) -- (0.561,0.647);
\draw[loop, line cap=round] (0.525,0.242) -- (0.499,0.252);
\draw[loop, line cap=round] (-0.499,0.595) -- (-0.525,0.583);
\draw[loop, line cap=round] (-0.561,0.304) -- (-0.594,0.296);
\filldraw[fill=slab!28!paper, draw=slab!28!paper, line width=0.2pt] (0.511,0.034) -- (0.491,0.045) -- (0.628,0.834) -- (0.664,0.823) -- cycle;
\filldraw[fill=slab!42!paper, draw=slab!42!paper, line width=0.2pt] (-0.628,0.153) -- (-0.664,0.146) -- (-0.511,0.712) -- (-0.491,0.726) -- cycle;
\draw[ink2, line width=0.6pt, line cap=round] (0.511,0.034) -- (0.491,0.045);
\draw[ink2, line width=0.6pt, line cap=round] (0.664,0.823) -- (0.628,0.834);
\draw[ink2, line width=0.6pt, line cap=round] (-0.628,0.153) -- (-0.664,0.146);
\draw[ink2, line width=0.6pt, line cap=round] (-0.491,0.726) -- (-0.511,0.712);
\draw[loop, line cap=round] (0.626,0.626) -- (0.594,0.637);
\draw[loop, line cap=round] (0.550,0.232) -- (0.525,0.242);
\draw[loop, line cap=round] (-0.525,0.583) -- (-0.550,0.570);
\draw[loop, line cap=round] (-0.594,0.296) -- (-0.626,0.287);
\filldraw[fill=slab!29!paper, draw=slab!29!paper, line width=0.2pt] (0.531,0.024) -- (0.511,0.034) -- (0.664,0.823) -- (0.700,0.811) -- cycle;
\filldraw[fill=slab!44!paper, draw=slab!44!paper, line width=0.2pt] (-0.664,0.146) -- (-0.700,0.138) -- (-0.531,0.697) -- (-0.511,0.712) -- cycle;
\draw[ink2, line width=0.6pt, line cap=round] (0.531,0.024) -- (0.511,0.034);
\draw[ink2, line width=0.6pt, line cap=round] (0.700,0.811) -- (0.664,0.823);
\draw[ink2, line width=0.6pt, line cap=round] (-0.664,0.146) -- (-0.700,0.138);
\draw[ink2, line width=0.6pt, line cap=round] (-0.511,0.712) -- (-0.531,0.697);
\draw[loop, line cap=round] (0.658,0.615) -- (0.626,0.626);
\draw[loop, line cap=round] (0.574,0.221) -- (0.550,0.232);
\draw[loop, line cap=round] (-0.550,0.570) -- (-0.574,0.557);
\draw[loop, line cap=round] (-0.626,0.287) -- (-0.658,0.278);
\filldraw[fill=slab!45!paper, draw=slab!45!paper, line width=0.2pt] (-0.700,0.138) -- (-0.735,0.130) -- (-0.550,0.682) -- (-0.531,0.697) -- cycle;
\draw[ink2, line width=0.6pt, line cap=round] (-0.700,0.138) -- (-0.735,0.130);
\draw[ink2, line width=0.6pt, line cap=round] (-0.531,0.697) -- (-0.550,0.682);
\draw[loop, line cap=round] (-0.574,0.557) -- (-0.597,0.544);
\filldraw[fill=slab!29!paper, draw=slab!29!paper, line width=0.2pt] (0.550,0.013) -- (0.531,0.024) -- (0.700,0.811) -- (0.735,0.799) -- cycle;
\draw[ink2, line width=0.6pt, line cap=round] (0.550,0.013) -- (0.531,0.024);
\draw[ink2, line width=0.6pt, line cap=round] (0.735,0.799) -- (0.700,0.811);
\draw[loop, line cap=round] (0.689,0.602) -- (0.658,0.615);
\draw[loop, line cap=round] (0.597,0.209) -- (0.574,0.221);
\filldraw[fill=slab!30!paper, draw=slab!30!paper, line width=0.2pt] (0.568,0.002) -- (0.550,0.013) -- (0.735,0.799) -- (0.770,0.786) -- cycle;
\filldraw[fill=slab!46!paper, draw=slab!46!paper, line width=0.2pt] (-0.735,0.130) -- (-0.770,0.121) -- (-0.568,0.667) -- (-0.550,0.682) -- cycle;
\draw[ink2, line width=0.6pt, line cap=round] (0.568,0.002) -- (0.550,0.013);
\draw[ink2, line width=0.6pt, line cap=round] (0.770,0.786) -- (0.735,0.799);
\draw[ink2, line width=0.6pt, line cap=round] (-0.735,0.130) -- (-0.770,0.121);
\draw[ink2, line width=0.6pt, line cap=round] (-0.550,0.682) -- (-0.568,0.667);
\draw[loop, line cap=round] (0.719,0.590) -- (0.689,0.602);
\draw[loop, line cap=round] (-0.597,0.544) -- (-0.619,0.530);
\filldraw[fill=slab!48!paper, draw=slab!48!paper, line width=0.2pt] (-0.770,0.121) -- (-0.804,0.111) -- (-0.585,0.651) -- (-0.568,0.667) -- cycle;
\draw[ink2, line width=0.6pt, line cap=round] (-0.770,0.121) -- (-0.804,0.111);
\draw[ink2, line width=0.6pt, line cap=round] (-0.568,0.667) -- (-0.585,0.651);
\draw[loop, line cap=round] (-0.619,0.530) -- (-0.640,0.516);
\filldraw[fill=slab!30!paper, draw=slab!30!paper, line width=0.2pt] (0.585,-0.009) -- (0.568,0.002) -- (0.770,0.786) -- (0.804,0.771) -- cycle;
\draw[ink2, line width=0.6pt, line cap=round] (0.585,-0.009) -- (0.568,0.002);
\draw[ink2, line width=0.6pt, line cap=round] (0.804,0.771) -- (0.770,0.786);
\draw[loop, line cap=round] (0.749,0.576) -- (0.719,0.590);
\filldraw[fill=slab!31!paper, draw=slab!31!paper, line width=0.2pt] (0.601,-0.020) -- (0.585,-0.009) -- (0.804,0.771) -- (0.837,0.756) -- cycle;
\draw[ink2, line width=0.6pt, line cap=round] (0.601,-0.020) -- (0.585,-0.009);
\draw[ink2, line width=0.6pt, line cap=round] (0.837,0.756) -- (0.804,0.771);
\draw[loop, line cap=round] (0.778,0.562) -- (0.749,0.576);
\filldraw[fill=slab!49!paper, draw=slab!49!paper, line width=0.2pt] (-0.804,0.111) -- (-0.837,0.101) -- (-0.601,0.636) -- (-0.585,0.651) -- cycle;
\draw[ink2, line width=0.6pt, line cap=round] (-0.804,0.111) -- (-0.837,0.101);
\draw[ink2, line width=0.6pt, line cap=round] (-0.585,0.651) -- (-0.601,0.636);
\draw[loop, line cap=round] (-0.640,0.516) -- (-0.660,0.502);
\draw[loop, line cap=round] (-0.749,0.246) -- (-0.778,0.234);
\filldraw[fill=slab!32!paper, draw=slab!32!paper, line width=0.2pt] (0.616,-0.031) -- (0.601,-0.020) -- (0.837,0.756) -- (0.870,0.741) -- cycle;
\draw[ink2, line width=0.6pt, line cap=round] (0.616,-0.031) -- (0.601,-0.020);
\draw[ink2, line width=0.6pt, line cap=round] (0.870,0.741) -- (0.837,0.756);
\draw[loop, line cap=round] (0.807,0.548) -- (0.778,0.562);
\draw[loop, line cap=round] (0.680,0.162) -- (0.660,0.174);
\filldraw[fill=slab!51!paper, draw=slab!51!paper, line width=0.2pt] (-0.837,0.101) -- (-0.870,0.090) -- (-0.616,0.619) -- (-0.601,0.636) -- cycle;
\draw[ink2, line width=0.6pt, line cap=round] (-0.837,0.101) -- (-0.870,0.090);
\draw[ink2, line width=0.6pt, line cap=round] (-0.601,0.636) -- (-0.616,0.619);
\draw[loop, line cap=round] (-0.660,0.502) -- (-0.680,0.487);
\draw[loop, line cap=round] (-0.778,0.234) -- (-0.807,0.222);
\filldraw[fill=slab!32!paper, draw=slab!32!paper, line width=0.2pt] (0.630,-0.043) -- (0.616,-0.031) -- (0.870,0.741) -- (0.902,0.724) -- cycle;
\filldraw[fill=slab!52!paper, draw=slab!52!paper, line width=0.2pt] (-0.870,0.090) -- (-0.902,0.078) -- (-0.630,0.603) -- (-0.616,0.619) -- cycle;
\draw[ink2, line width=0.6pt, line cap=round] (0.630,-0.043) -- (0.616,-0.031);
\draw[ink2, line width=0.6pt, line cap=round] (0.902,0.724) -- (0.870,0.741);
\draw[ink2, line width=0.6pt, line cap=round] (-0.870,0.090) -- (-0.902,0.078);
\draw[ink2, line width=0.6pt, line cap=round] (-0.616,0.619) -- (-0.630,0.603);
\draw[loop, line cap=round] (0.834,0.532) -- (0.807,0.548);
\draw[loop, line cap=round] (0.698,0.149) -- (0.680,0.162);
\draw[loop, line cap=round] (-0.680,0.487) -- (-0.698,0.472);
\draw[loop, line cap=round] (-0.807,0.222) -- (-0.834,0.209);
\filldraw[fill=slab!33!paper, draw=slab!33!paper, line width=0.2pt] (0.643,-0.054) -- (0.630,-0.043) -- (0.902,0.724) -- (0.933,0.706) -- cycle;
\filldraw[fill=slab!54!paper, draw=slab!54!paper, line width=0.2pt] (-0.902,0.078) -- (-0.933,0.066) -- (-0.643,0.586) -- (-0.630,0.603) -- cycle;
\draw[ink2, line width=0.6pt, line cap=round] (0.643,-0.054) -- (0.630,-0.043);
\draw[ink2, line width=0.6pt, line cap=round] (0.933,0.706) -- (0.902,0.724);
\draw[ink2, line width=0.6pt, line cap=round] (-0.902,0.078) -- (-0.933,0.066);
\draw[ink2, line width=0.6pt, line cap=round] (-0.630,0.603) -- (-0.643,0.586);
\draw[loop, line cap=round] (0.861,0.516) -- (0.834,0.532);
\draw[loop, line cap=round] (0.715,0.136) -- (0.698,0.149);
\draw[loop, line cap=round] (-0.698,0.472) -- (-0.715,0.456);
\draw[loop, line cap=round] (-0.834,0.209) -- (-0.861,0.196);
\filldraw[fill=slab!34!paper, draw=slab!34!paper, line width=0.2pt] (0.655,-0.065) -- (0.643,-0.054) -- (0.933,0.706) -- (0.963,0.688) -- cycle;
\filldraw[fill=slab!55!paper, draw=slab!55!paper, line width=0.2pt] (-0.933,0.066) -- (-0.963,0.054) -- (-0.655,0.569) -- (-0.643,0.586) -- cycle;
\draw[ink2, line width=0.6pt, line cap=round] (0.655,-0.065) -- (0.643,-0.054);
\draw[ink2, line width=0.6pt, line cap=round] (0.963,0.688) -- (0.933,0.706);
\draw[ink2, line width=0.6pt, line cap=round] (-0.933,0.066) -- (-0.963,0.054);
\draw[ink2, line width=0.6pt, line cap=round] (-0.643,0.586) -- (-0.655,0.569);
\draw[loop, line cap=round] (0.886,0.500) -- (0.861,0.516);
\draw[loop, line cap=round] (0.732,0.123) -- (0.715,0.136);
\draw[loop, line cap=round] (-0.715,0.456) -- (-0.732,0.440);
\filldraw[fill=slab!35!paper, draw=slab!35!paper, line width=0.2pt] (0.666,-0.077) -- (0.655,-0.065) -- (0.963,0.688) -- (0.993,0.669) -- cycle;
\draw[ink2, line width=0.6pt, line cap=round] (0.666,-0.077) -- (0.655,-0.065);
\draw[ink2, line width=0.6pt, line cap=round] (0.993,0.669) -- (0.963,0.688);
\draw[loop, line cap=round] (0.911,0.483) -- (0.886,0.500);
\filldraw[fill=slab!57!paper, draw=slab!57!paper, line width=0.2pt] (-0.963,0.054) -- (-0.993,0.040) -- (-0.666,0.552) -- (-0.655,0.569) -- cycle;
\draw[ink2, line width=0.6pt, line cap=round] (-0.963,0.054) -- (-0.993,0.040);
\draw[ink2, line width=0.6pt, line cap=round] (-0.655,0.569) -- (-0.666,0.552);
\draw[loop, line cap=round] (-0.732,0.440) -- (-0.747,0.424);
\filldraw[fill=slab!35!paper, draw=slab!35!paper, line width=0.2pt] (0.675,-0.088) -- (0.666,-0.077) -- (0.993,0.669) -- (1.021,0.650) -- cycle;
\draw[ink2, line width=0.6pt, line cap=round] (0.675,-0.088) -- (0.666,-0.077);
\draw[ink2, line width=0.6pt, line cap=round] (1.021,0.650) -- (0.993,0.669);
\draw[loop, line cap=round] (0.934,0.465) -- (0.911,0.483);
\filldraw[fill=slab!58!paper, draw=slab!58!paper, line width=0.2pt] (-0.993,0.040) -- (-1.021,0.026) -- (-0.675,0.535) -- (-0.666,0.552) -- cycle;
\draw[ink2, line width=0.6pt, line cap=round] (-0.993,0.040) -- (-1.021,0.026);
\draw[ink2, line width=0.6pt, line cap=round] (-0.666,0.552) -- (-0.675,0.535);
\draw[loop, line cap=round] (-0.747,0.424) -- (-0.762,0.408);
\filldraw[fill=slab!60!paper, draw=slab!60!paper, line width=0.2pt] (-1.021,0.026) -- (-1.048,0.012) -- (-0.684,0.518) -- (-0.675,0.535) -- cycle;
\draw[ink2, line width=0.6pt, line cap=round] (-1.021,0.026) -- (-1.048,0.012);
\draw[ink2, line width=0.6pt, line cap=round] (-0.675,0.535) -- (-0.684,0.518);
\draw[loop, line cap=round] (-0.762,0.408) -- (-0.775,0.391);
\draw[loop, line cap=round] (-0.934,0.154) -- (-0.957,0.139);
\filldraw[fill=slab!36!paper, draw=slab!36!paper, line width=0.2pt] (0.684,-0.099) -- (0.675,-0.088) -- (1.021,0.650) -- (1.048,0.629) -- cycle;
\draw[ink2, line width=0.6pt, line cap=round] (0.684,-0.099) -- (0.675,-0.088);
\draw[ink2, line width=0.6pt, line cap=round] (1.048,0.629) -- (1.021,0.650);
\draw[ink3!50, line width=0.25pt] (0.684,-0.099) -- (1.048,0.629);
\draw[loop, line cap=round] (0.957,0.447) -- (0.934,0.465);
\filldraw[fill=slab!61!paper, draw=slab!61!paper, line width=0.2pt] (-1.048,0.012) -- (-1.074,-0.003) -- (-0.692,0.500) -- (-0.684,0.518) -- cycle;
\draw[ink2, line width=0.6pt, line cap=round] (-1.048,0.012) -- (-1.074,-0.003);
\draw[ink2, line width=0.6pt, line cap=round] (-0.684,0.518) -- (-0.692,0.500);
\draw[ink3!50, line width=0.25pt] (-1.048,0.012) -- (-0.684,0.518);
\draw[loop, line cap=round] (-0.775,0.391) -- (-0.787,0.375);
\draw[loop, line cap=round] (-0.957,0.139) -- (-0.978,0.123);
\filldraw[fill=slab!37!paper, draw=slab!37!paper, line width=0.2pt] (0.692,-0.110) -- (0.684,-0.099) -- (1.048,0.629) -- (1.074,0.608) -- cycle;
\draw[ink2, line width=0.6pt, line cap=round] (0.692,-0.110) -- (0.684,-0.099);
\draw[ink2, line width=0.6pt, line cap=round] (1.074,0.608) -- (1.048,0.629);
\draw[loop, line cap=round] (0.978,0.428) -- (0.957,0.447);
\draw[loop, line cap=round] (0.787,0.069) -- (0.775,0.083);
\filldraw[fill=slab!38!paper, draw=slab!38!paper, line width=0.2pt] (0.699,-0.121) -- (0.692,-0.110) -- (1.074,0.608) -- (1.099,0.586) -- cycle;
\draw[ink2, line width=0.6pt, line cap=round] (0.699,-0.121) -- (0.692,-0.110);
\draw[ink2, line width=0.6pt, line cap=round] (1.099,0.586) -- (1.074,0.608);
\draw[loop, line cap=round] (0.999,0.409) -- (0.978,0.428);
\draw[loop, line cap=round] (0.799,0.055) -- (0.787,0.069);
\filldraw[fill=slab!63!paper, draw=slab!63!paper, line width=0.2pt] (-1.074,-0.003) -- (-1.099,-0.018) -- (-0.699,0.483) -- (-0.692,0.500) -- cycle;
\draw[ink2, line width=0.6pt, line cap=round] (-1.074,-0.003) -- (-1.099,-0.018);
\draw[ink2, line width=0.6pt, line cap=round] (-0.692,0.500) -- (-0.699,0.483);
\draw[loop, line cap=round] (-0.787,0.375) -- (-0.799,0.357);
\draw[loop, line cap=round] (-0.978,0.123) -- (-0.999,0.107);
\filldraw[fill=slab!39!paper, draw=slab!39!paper, line width=0.2pt] (0.705,-0.132) -- (0.699,-0.121) -- (1.099,0.586) -- (1.123,0.564) -- cycle;
\draw[ink2, line width=0.6pt, line cap=round] (0.705,-0.132) -- (0.699,-0.121);
\draw[ink2, line width=0.6pt, line cap=round] (1.123,0.564) -- (1.099,0.586);
\draw[loop, line cap=round] (1.018,0.390) -- (0.999,0.409);
\draw[loop, line cap=round] (0.809,0.042) -- (0.799,0.055);
\filldraw[fill=slab!64!paper, draw=slab!64!paper, line width=0.2pt] (-1.099,-0.018) -- (-1.123,-0.034) -- (-0.705,0.465) -- (-0.699,0.483) -- cycle;
\draw[ink2, line width=0.6pt, line cap=round] (-1.099,-0.018) -- (-1.123,-0.034);
\draw[ink2, line width=0.6pt, line cap=round] (-0.699,0.483) -- (-0.705,0.465);
\draw[loop, line cap=round] (-0.799,0.357) -- (-0.809,0.340);
\draw[loop, line cap=round] (-0.999,0.107) -- (-1.018,0.091);
\filldraw[fill=slab!39!paper, draw=slab!39!paper, line width=0.2pt] (0.709,-0.143) -- (0.705,-0.132) -- (1.123,0.564) -- (1.145,0.540) -- cycle;
\draw[ink2, line width=0.6pt, line cap=round] (0.709,-0.143) -- (0.705,-0.132);
\draw[ink2, line width=0.6pt, line cap=round] (1.145,0.540) -- (1.123,0.564);
\draw[loop, line cap=round] (1.036,0.369) -- (1.018,0.390);
\draw[loop, line cap=round] (0.818,0.028) -- (0.809,0.042);
\filldraw[fill=slab!64!paper, draw=slab!64!paper, line width=0.2pt] (-1.123,-0.034) -- (-1.145,-0.050) -- (-0.709,0.447) -- (-0.705,0.465) -- cycle;
\draw[ink2, line width=0.6pt, line cap=round] (-1.123,-0.034) -- (-1.145,-0.050);
\draw[ink2, line width=0.6pt, line cap=round] (-0.705,0.465) -- (-0.709,0.447);
\draw[loop, line cap=round] (-0.809,0.340) -- (-0.818,0.323);
\filldraw[fill=slab!40!paper, draw=slab!40!paper, line width=0.2pt] (0.713,-0.154) -- (0.709,-0.143) -- (1.145,0.540) -- (1.166,0.517) -- cycle;
\draw[ink2, line width=0.6pt, line cap=round] (0.713,-0.154) -- (0.709,-0.143);
\draw[ink2, line width=0.6pt, line cap=round] (1.166,0.517) -- (1.145,0.540);
\draw[loop, line cap=round] (1.053,0.349) -- (1.036,0.369);
\filldraw[fill=slab!63!paper, draw=slab!63!paper, line width=0.2pt] (-1.145,-0.050) -- (-1.166,-0.067) -- (-0.713,0.429) -- (-0.709,0.447) -- cycle;
\draw[ink2, line width=0.6pt, line cap=round] (-1.145,-0.050) -- (-1.166,-0.067);
\draw[ink2, line width=0.6pt, line cap=round] (-0.709,0.447) -- (-0.713,0.429);
\draw[loop, line cap=round] (-0.818,0.323) -- (-0.827,0.305);
\filldraw[fill=slab!41!paper, draw=slab!41!paper, line width=0.2pt] (0.716,-0.165) -- (0.713,-0.154) -- (1.166,0.517) -- (1.186,0.492) -- cycle;
\draw[ink2, line width=0.6pt, line cap=round] (0.716,-0.165) -- (0.713,-0.154);
\draw[ink2, line width=0.6pt, line cap=round] (1.186,0.492) -- (1.166,0.517);
\draw[loop, line cap=round] (1.068,0.328) -- (1.053,0.349);
\filldraw[fill=slab!61!paper, draw=slab!61!paper, line width=0.2pt] (-1.166,-0.067) -- (-1.186,-0.084) -- (-0.716,0.411) -- (-0.713,0.429) -- cycle;
\draw[ink2, line width=0.6pt, line cap=round] (-1.166,-0.067) -- (-1.186,-0.084);
\draw[ink2, line width=0.6pt, line cap=round] (-0.713,0.429) -- (-0.716,0.411);
\draw[loop, line cap=round] (-0.827,0.305) -- (-0.834,0.288);
\filldraw[fill=slab!42!paper, draw=slab!42!paper, line width=0.2pt] (0.718,-0.175) -- (0.716,-0.165) -- (1.186,0.492) -- (1.204,0.467) -- cycle;
\draw[ink2, line width=0.6pt, line cap=round] (0.718,-0.175) -- (0.716,-0.165);
\draw[ink2, line width=0.6pt, line cap=round] (1.204,0.467) -- (1.186,0.492);
\draw[loop, line cap=round] (1.083,0.307) -- (1.068,0.328);
\filldraw[fill=slab!60!paper, draw=slab!60!paper, line width=0.2pt] (-1.186,-0.084) -- (-1.204,-0.101) -- (-0.718,0.394) -- (-0.716,0.411) -- cycle;
\draw[ink2, line width=0.6pt, line cap=round] (-1.186,-0.084) -- (-1.204,-0.101);
\draw[ink2, line width=0.6pt, line cap=round] (-0.716,0.411) -- (-0.718,0.394);
\draw[loop, line cap=round] (-0.834,0.288) -- (-0.840,0.270);
\draw[loop, line cap=round] (-1.068,0.040) -- (-1.083,0.022);
\filldraw[fill=slab!43!paper, draw=slab!43!paper, line width=0.2pt] (0.719,-0.186) -- (0.718,-0.175) -- (1.204,0.467) -- (1.221,0.442) -- cycle;
\draw[ink2, line width=0.6pt, line cap=round] (0.719,-0.186) -- (0.718,-0.175);
\draw[ink2, line width=0.6pt, line cap=round] (1.221,0.442) -- (1.204,0.467);
\draw[loop, line cap=round] (1.096,0.285) -- (1.083,0.307);
\draw[loop, line cap=round] (0.845,-0.029) -- (0.840,-0.015);
\filldraw[fill=slab!58!paper, draw=slab!58!paper, line width=0.2pt] (-1.204,-0.101) -- (-1.221,-0.119) -- (-0.719,0.376) -- (-0.718,0.394) -- cycle;
\draw[ink2, line width=0.6pt, line cap=round] (-1.204,-0.101) -- (-1.221,-0.119);
\draw[ink2, line width=0.6pt, line cap=round] (-0.718,0.394) -- (-0.719,0.376);
\draw[loop, line cap=round] (-0.840,0.270) -- (-0.845,0.252);
\draw[loop, line cap=round] (-1.083,0.022) -- (-1.096,0.004);
\filldraw[fill=slab!57!paper, draw=slab!57!paper, line width=0.2pt] (-1.221,-0.119) -- (-1.237,-0.138) -- (-0.720,0.358) -- (-0.719,0.376) -- cycle;
\draw[ink2, line width=0.6pt, line cap=round] (-1.221,-0.119) -- (-1.237,-0.138);
\draw[ink2, line width=0.6pt, line cap=round] (-0.719,0.376) -- (-0.720,0.358);
\draw[loop, line cap=round] (-0.845,0.252) -- (-0.849,0.234);
\draw[loop, line cap=round] (-1.096,0.004) -- (-1.107,-0.014);
\filldraw[fill=slab!43!paper, draw=slab!43!paper, line width=0.2pt] (0.720,-0.196) -- (0.719,-0.186) -- (1.221,0.442) -- (1.237,0.416) -- cycle;
\draw[ink2, line width=0.6pt, line cap=round] (0.720,-0.196) -- (0.719,-0.186);
\draw[ink2, line width=0.6pt, line cap=round] (1.237,0.416) -- (1.221,0.442);
\draw[loop, line cap=round] (1.107,0.263) -- (1.096,0.285);
\draw[loop, line cap=round] (0.849,-0.043) -- (0.845,-0.029);
\filldraw[fill=slab!55!paper, draw=slab!55!paper, line width=0.2pt] (-1.237,-0.138) -- (-1.251,-0.156) -- (-0.719,0.340) -- (-0.720,0.358) -- cycle;
\draw[ink2, line width=0.6pt, line cap=round] (-1.237,-0.138) -- (-1.251,-0.156);
\draw[ink2, line width=0.6pt, line cap=round] (-0.720,0.358) -- (-0.719,0.340);
\draw[loop, line cap=round] (-0.849,0.234) -- (-0.852,0.216);
\draw[loop, line cap=round] (-1.107,-0.014) -- (-1.118,-0.032);
\filldraw[fill=slab!44!paper, draw=slab!44!paper, line width=0.2pt] (0.719,-0.206) -- (0.720,-0.196) -- (1.237,0.416) -- (1.251,0.390) -- cycle;
\draw[ink2, line width=0.6pt, line cap=round] (0.719,-0.206) -- (0.720,-0.196);
\draw[ink2, line width=0.6pt, line cap=round] (1.251,0.390) -- (1.237,0.416);
\draw[loop, line cap=round] (1.118,0.241) -- (1.107,0.263);
\draw[loop, line cap=round] (0.852,-0.057) -- (0.849,-0.043);
\filldraw[fill=slab!54!paper, draw=slab!54!paper, line width=0.2pt] (-1.251,-0.156) -- (-1.263,-0.175) -- (-0.717,0.322) -- (-0.719,0.340) -- cycle;
\draw[ink2, line width=0.6pt, line cap=round] (-1.251,-0.156) -- (-1.263,-0.175);
\draw[ink2, line width=0.6pt, line cap=round] (-0.719,0.340) -- (-0.717,0.322);
\draw[loop, line cap=round] (-0.852,0.216) -- (-0.854,0.198);
\filldraw[fill=slab!45!paper, draw=slab!45!paper, line width=0.2pt] (0.717,-0.215) -- (0.719,-0.206) -- (1.251,0.390) -- (1.263,0.363) -- cycle;
\draw[ink2, line width=0.6pt, line cap=round] (0.717,-0.215) -- (0.719,-0.206);
\draw[ink2, line width=0.6pt, line cap=round] (1.263,0.363) -- (1.251,0.390);
\draw[loop, line cap=round] (1.127,0.218) -- (1.118,0.241);
\draw[loop, line cap=round] (0.854,-0.071) -- (0.852,-0.057);
\filldraw[fill=slab!46!paper, draw=slab!46!paper, line width=0.2pt] (0.715,-0.225) -- (0.717,-0.215) -- (1.263,0.363) -- (1.274,0.336) -- cycle;
\filldraw[fill=slab!53!paper, draw=slab!53!paper, line width=0.2pt] (-1.263,-0.175) -- (-1.274,-0.194) -- (-0.715,0.305) -- (-0.717,0.322) -- cycle;
\draw[ink2, line width=0.6pt, line cap=round] (0.715,-0.225) -- (0.717,-0.215);
\draw[ink2, line width=0.6pt, line cap=round] (1.274,0.336) -- (1.263,0.363);
\draw[ink2, line width=0.6pt, line cap=round] (-1.263,-0.175) -- (-1.274,-0.194);
\draw[ink2, line width=0.6pt, line cap=round] (-0.717,0.322) -- (-0.715,0.305);
\draw[loop, line cap=round] (1.134,0.196) -- (1.127,0.218);
\draw[loop, line cap=round] (-0.854,0.198) -- (-0.855,0.180);
\filldraw[fill=slab!46!paper, draw=slab!46!paper, line width=0.2pt] (0.712,-0.234) -- (0.715,-0.225) -- (1.274,0.336) -- (1.283,0.308) -- cycle;
\draw[ink2, line width=0.6pt, line cap=round] (0.712,-0.234) -- (0.715,-0.225);
\draw[ink2, line width=0.6pt, line cap=round] (1.283,0.308) -- (1.274,0.336);
\draw[loop, line cap=round] (1.140,0.173) -- (1.134,0.196);
\filldraw[fill=slab!51!paper, draw=slab!51!paper, line width=0.2pt] (-1.274,-0.194) -- (-1.283,-0.213) -- (-0.712,0.287) -- (-0.715,0.305) -- cycle;
\draw[ink2, line width=0.6pt, line cap=round] (-1.274,-0.194) -- (-1.283,-0.213);
\draw[ink2, line width=0.6pt, line cap=round] (-0.715,0.305) -- (-0.712,0.287);
\draw[loop, line cap=round] (-0.855,0.180) -- (-0.855,0.162);
\filldraw[fill=slab!47!paper, draw=slab!47!paper, line width=0.2pt] (0.708,-0.243) -- (0.712,-0.234) -- (1.283,0.308) -- (1.291,0.280) -- cycle;
\draw[ink2, line width=0.6pt, line cap=round] (0.708,-0.243) -- (0.712,-0.234);
\draw[ink2, line width=0.6pt, line cap=round] (1.291,0.280) -- (1.283,0.308);
\draw[loop, line cap=round] (1.145,0.149) -- (1.140,0.173);
\filldraw[fill=slab!50!paper, draw=slab!50!paper, line width=0.2pt] (-1.283,-0.213) -- (-1.291,-0.232) -- (-0.708,0.269) -- (-0.712,0.287) -- cycle;
\draw[ink2, line width=0.6pt, line cap=round] (-1.283,-0.213) -- (-1.291,-0.232);
\draw[ink2, line width=0.6pt, line cap=round] (-0.712,0.287) -- (-0.708,0.269);
\draw[loop, line cap=round] (-0.855,0.162) -- (-0.854,0.144);
\draw[loop, line cap=round] (-1.140,-0.088) -- (-1.145,-0.107);
\filldraw[fill=slab!48!paper, draw=slab!48!paper, line width=0.2pt] (0.703,-0.252) -- (0.708,-0.243) -- (1.291,0.280) -- (1.297,0.252) -- cycle;
\draw[ink2, line width=0.6pt, line cap=round] (0.703,-0.252) -- (0.708,-0.243);
\draw[ink2, line width=0.6pt, line cap=round] (1.297,0.252) -- (1.291,0.280);
\draw[ink3!50, line width=0.25pt] (0.703,-0.252) -- (1.297,0.252);
\draw[loop, line cap=round] (1.148,0.126) -- (1.145,0.149);
\draw[loop, line cap=round] (0.852,-0.126) -- (0.854,-0.112);
\filldraw[fill=slab!49!paper, draw=slab!49!paper, line width=0.2pt] (-1.291,-0.232) -- (-1.297,-0.252) -- (-0.703,0.252) -- (-0.708,0.269) -- cycle;
\draw[ink2, line width=0.6pt, line cap=round] (-1.291,-0.232) -- (-1.297,-0.252);
\draw[ink2, line width=0.6pt, line cap=round] (-0.708,0.269) -- (-0.703,0.252);
\draw[loop, line cap=round] (-0.854,0.144) -- (-0.852,0.126);
\draw[loop, line cap=round] (-1.145,-0.107) -- (-1.148,-0.126);
\filldraw[fill=slab!48!paper, draw=slab!48!paper, line width=0.2pt] (0.697,-0.260) -- (0.703,-0.252) -- (1.297,0.252) -- (1.301,0.223) -- cycle;
\draw[ink2, line width=0.6pt, line cap=round] (0.697,-0.260) -- (0.703,-0.252);
\draw[ink2, line width=0.6pt, line cap=round] (1.301,0.223) -- (1.297,0.252);
\draw[loop, line cap=round] (1.150,0.102) -- (1.148,0.126);
\draw[loop, line cap=round] (0.848,-0.139) -- (0.852,-0.126);
\filldraw[fill=slab!48!paper, draw=slab!48!paper, line width=0.2pt] (-1.297,-0.252) -- (-1.301,-0.272) -- (-0.697,0.235) -- (-0.703,0.252) -- cycle;
\draw[ink2, line width=0.6pt, line cap=round] (-1.297,-0.252) -- (-1.301,-0.272);
\draw[ink2, line width=0.6pt, line cap=round] (-0.703,0.252) -- (-0.697,0.235);
\draw[ink3!50, line width=0.25pt] (-1.297,-0.252) -- (-0.703,0.252);
\draw[loop, line cap=round] (-0.852,0.126) -- (-0.848,0.108);
\draw[loop, line cap=round] (-1.148,-0.126) -- (-1.150,-0.145);
\filldraw[fill=slab!49!paper, draw=slab!49!paper, line width=0.2pt] (0.691,-0.269) -- (0.697,-0.260) -- (1.301,0.223) -- (1.304,0.195) -- cycle;
\draw[ink2, line width=0.6pt, line cap=round] (0.691,-0.269) -- (0.697,-0.260);
\draw[ink2, line width=0.6pt, line cap=round] (1.304,0.195) -- (1.301,0.223);
\draw[loop, line cap=round] (1.151,0.079) -- (1.150,0.102);
\draw[loop, line cap=round] (0.844,-0.153) -- (0.848,-0.139);
\filldraw[fill=slab!46!paper, draw=slab!46!paper, line width=0.2pt] (-1.301,-0.272) -- (-1.304,-0.291) -- (-0.691,0.217) -- (-0.697,0.235) -- cycle;
\draw[ink2, line width=0.6pt, line cap=round] (-1.301,-0.272) -- (-1.304,-0.291);
\draw[ink2, line width=0.6pt, line cap=round] (-0.697,0.235) -- (-0.691,0.217);
\draw[loop, line cap=round] (-0.848,0.108) -- (-0.844,0.090);
\draw[loop, line cap=round] (-1.150,-0.145) -- (-1.151,-0.164);
\filldraw[fill=slab!45!paper, draw=slab!45!paper, line width=0.2pt] (-1.304,-0.291) -- (-1.305,-0.311) -- (-0.684,0.200) -- (-0.691,0.217) -- cycle;
\draw[ink2, line width=0.6pt, line cap=round] (-1.304,-0.291) -- (-1.305,-0.311);
\draw[ink2, line width=0.6pt, line cap=round] (-0.691,0.217) -- (-0.684,0.200);
\draw[loop, line cap=round] (-0.844,0.090) -- (-0.839,0.072);
\filldraw[fill=slab!50!paper, draw=slab!50!paper, line width=0.2pt] (0.684,-0.276) -- (0.691,-0.269) -- (1.304,0.195) -- (1.305,0.166) -- cycle;
\draw[ink2, line width=0.6pt, line cap=round] (0.684,-0.276) -- (0.691,-0.269);
\draw[ink2, line width=0.6pt, line cap=round] (1.305,0.166) -- (1.304,0.195);
\draw[loop, line cap=round] (1.150,0.055) -- (1.151,0.079);
\draw[loop, line cap=round] (0.839,-0.166) -- (0.844,-0.153);
\filldraw[fill=slab!44!paper, draw=slab!44!paper, line width=0.2pt] (-1.305,-0.311) -- (-1.304,-0.331) -- (-0.676,0.183) -- (-0.684,0.200) -- cycle;
\draw[ink2, line width=0.6pt, line cap=round] (-1.305,-0.311) -- (-1.304,-0.331);
\draw[ink2, line width=0.6pt, line cap=round] (-0.684,0.200) -- (-0.676,0.183);
\draw[loop, line cap=round] (-0.839,0.072) -- (-0.833,0.055);
\filldraw[fill=slab!50!paper, draw=slab!50!paper, line width=0.2pt] (0.676,-0.284) -- (0.684,-0.276) -- (1.305,0.166) -- (1.304,0.137) -- cycle;
\draw[ink2, line width=0.6pt, line cap=round] (0.676,-0.284) -- (0.684,-0.276);
\draw[ink2, line width=0.6pt, line cap=round] (1.304,0.137) -- (1.305,0.166);
\draw[loop, line cap=round] (1.147,0.031) -- (1.150,0.055);
\filldraw[fill=slab!43!paper, draw=slab!43!paper, line width=0.2pt] (-1.304,-0.331) -- (-1.302,-0.351) -- (-0.668,0.167) -- (-0.676,0.183) -- cycle;
\draw[ink2, line width=0.6pt, line cap=round] (-1.304,-0.331) -- (-1.302,-0.351);
\draw[ink2, line width=0.6pt, line cap=round] (-0.676,0.183) -- (-0.668,0.167);
\draw[loop, line cap=round] (-0.833,0.055) -- (-0.826,0.037);
\filldraw[fill=slab!51!paper, draw=slab!51!paper, line width=0.2pt] (0.668,-0.291) -- (0.676,-0.284) -- (1.304,0.137) -- (1.302,0.107) -- cycle;
\draw[ink2, line width=0.6pt, line cap=round] (0.668,-0.291) -- (0.676,-0.284);
\draw[ink2, line width=0.6pt, line cap=round] (1.302,0.107) -- (1.304,0.137);
\draw[loop, line cap=round] (1.143,0.008) -- (1.147,0.031);
\filldraw[fill=slab!42!paper, draw=slab!42!paper, line width=0.2pt] (-1.302,-0.351) -- (-1.297,-0.370) -- (-0.659,0.150) -- (-0.668,0.167) -- cycle;
\draw[ink2, line width=0.6pt, line cap=round] (-1.302,-0.351) -- (-1.297,-0.370);
\draw[ink2, line width=0.6pt, line cap=round] (-0.668,0.167) -- (-0.659,0.150);
\draw[loop, line cap=round] (-0.826,0.037) -- (-0.819,0.020);
\draw[loop, line cap=round] (-1.143,-0.221) -- (-1.138,-0.240);
\filldraw[fill=slab!51!paper, draw=slab!51!paper, line width=0.2pt] (0.659,-0.298) -- (0.668,-0.291) -- (1.302,0.107) -- (1.297,0.078) -- cycle;
\draw[ink2, line width=0.6pt, line cap=round] (0.659,-0.298) -- (0.668,-0.291);
\draw[ink2, line width=0.6pt, line cap=round] (1.297,0.078) -- (1.302,0.107);
\draw[loop, line cap=round] (1.138,-0.016) -- (1.143,0.008);
\filldraw[fill=slab!52!paper, draw=slab!52!paper, line width=0.2pt] (0.649,-0.305) -- (0.659,-0.298) -- (1.297,0.078) -- (1.291,0.049) -- cycle;
\filldraw[fill=slab!41!paper, draw=slab!41!paper, line width=0.2pt] (-1.297,-0.370) -- (-1.291,-0.390) -- (-0.649,0.134) -- (-0.659,0.150) -- cycle;
\draw[ink2, line width=0.6pt, line cap=round] (0.649,-0.305) -- (0.659,-0.298);
\draw[ink2, line width=0.6pt, line cap=round] (1.291,0.049) -- (1.297,0.078);
\draw[ink2, line width=0.6pt, line cap=round] (-1.297,-0.370) -- (-1.291,-0.390);
\draw[ink2, line width=0.6pt, line cap=round] (-0.659,0.150) -- (-0.649,0.134);
\draw[loop, line cap=round] (1.131,-0.040) -- (1.138,-0.016);
\draw[loop, line cap=round] (0.810,-0.217) -- (0.819,-0.204);
\draw[loop, line cap=round] (-0.819,0.020) -- (-0.810,0.003);
\draw[loop, line cap=round] (-1.138,-0.240) -- (-1.131,-0.259);
\filldraw[fill=slab!52!paper, draw=slab!52!paper, line width=0.2pt] (0.639,-0.311) -- (0.649,-0.305) -- (1.291,0.049) -- (1.284,0.019) -- cycle;
\draw[ink2, line width=0.6pt, line cap=round] (0.639,-0.311) -- (0.649,-0.305);
\draw[ink2, line width=0.6pt, line cap=round] (1.284,0.019) -- (1.291,0.049);
\draw[loop, line cap=round] (1.122,-0.063) -- (1.131,-0.040);
\draw[loop, line cap=round] (0.800,-0.229) -- (0.810,-0.217);
\filldraw[fill=slab!40!paper, draw=slab!40!paper, line width=0.2pt] (-1.291,-0.390) -- (-1.284,-0.409) -- (-0.639,0.117) -- (-0.649,0.134) -- cycle;
\draw[ink2, line width=0.6pt, line cap=round] (-1.291,-0.390) -- (-1.284,-0.409);
\draw[ink2, line width=0.6pt, line cap=round] (-0.649,0.134) -- (-0.639,0.117);
\draw[loop, line cap=round] (-0.810,0.003) -- (-0.800,-0.014);
\draw[loop, line cap=round] (-1.131,-0.259) -- (-1.122,-0.278);
\filldraw[fill=slab!52!paper, draw=slab!52!paper, line width=0.2pt] (0.628,-0.317) -- (0.639,-0.311) -- (1.284,0.019) -- (1.274,-0.010) -- cycle;
\draw[ink2, line width=0.6pt, line cap=round] (0.628,-0.317) -- (0.639,-0.311);
\draw[ink2, line width=0.6pt, line cap=round] (1.274,-0.010) -- (1.284,0.019);
\draw[loop, line cap=round] (1.113,-0.087) -- (1.122,-0.063);
\draw[loop, line cap=round] (0.789,-0.241) -- (0.800,-0.229);
\filldraw[fill=slab!39!paper, draw=slab!39!paper, line width=0.2pt] (-1.284,-0.409) -- (-1.274,-0.429) -- (-0.628,0.101) -- (-0.639,0.117) -- cycle;
\draw[ink2, line width=0.6pt, line cap=round] (-1.284,-0.409) -- (-1.274,-0.429);
\draw[ink2, line width=0.6pt, line cap=round] (-0.639,0.117) -- (-0.628,0.101);
\draw[loop, line cap=round] (-0.800,-0.014) -- (-0.789,-0.031);
\draw[loop, line cap=round] (-1.122,-0.278) -- (-1.113,-0.296);
\filldraw[fill=slab!53!paper, draw=slab!53!paper, line width=0.2pt] (0.616,-0.323) -- (0.628,-0.317) -- (1.274,-0.010) -- (1.263,-0.039) -- cycle;
\draw[ink2, line width=0.6pt, line cap=round] (0.616,-0.323) -- (0.628,-0.317);
\draw[ink2, line width=0.6pt, line cap=round] (1.263,-0.039) -- (1.274,-0.010);
\draw[loop, line cap=round] (1.101,-0.110) -- (1.113,-0.087);
\draw[loop, line cap=round] (0.778,-0.252) -- (0.789,-0.241);
\filldraw[fill=slab!38!paper, draw=slab!38!paper, line width=0.2pt] (-1.274,-0.429) -- (-1.263,-0.448) -- (-0.616,0.085) -- (-0.628,0.101) -- cycle;
\draw[ink2, line width=0.6pt, line cap=round] (-1.274,-0.429) -- (-1.263,-0.448);
\draw[ink2, line width=0.6pt, line cap=round] (-0.628,0.101) -- (-0.616,0.085);
\draw[loop, line cap=round] (-0.789,-0.031) -- (-0.778,-0.048);
\filldraw[fill=slab!53!paper, draw=slab!53!paper, line width=0.2pt] (0.604,-0.329) -- (0.616,-0.323) -- (1.263,-0.039) -- (1.250,-0.068) -- cycle;
\draw[ink2, line width=0.6pt, line cap=round] (0.604,-0.329) -- (0.616,-0.323);
\draw[ink2, line width=0.6pt, line cap=round] (1.250,-0.068) -- (1.263,-0.039);
\draw[loop, line cap=round] (1.089,-0.133) -- (1.101,-0.110);
\filldraw[fill=slab!38!paper, draw=slab!38!paper, line width=0.2pt] (-1.263,-0.448) -- (-1.250,-0.467) -- (-0.604,0.070) -- (-0.616,0.085) -- cycle;
\draw[ink2, line width=0.6pt, line cap=round] (-1.263,-0.448) -- (-1.250,-0.467);
\draw[ink2, line width=0.6pt, line cap=round] (-0.616,0.085) -- (-0.604,0.070);
\draw[loop, line cap=round] (-0.778,-0.048) -- (-0.766,-0.064);
\filldraw[fill=slab!37!paper, draw=slab!37!paper, line width=0.2pt] (-1.250,-0.467) -- (-1.235,-0.485) -- (-0.592,0.054) -- (-0.604,0.070) -- cycle;
\draw[ink2, line width=0.6pt, line cap=round] (-1.250,-0.467) -- (-1.235,-0.485);
\draw[ink2, line width=0.6pt, line cap=round] (-0.604,0.070) -- (-0.592,0.054);
\draw[loop, line cap=round] (-0.766,-0.064) -- (-0.753,-0.081);
\filldraw[fill=slab!53!paper, draw=slab!53!paper, line width=0.2pt] (0.592,-0.334) -- (0.604,-0.329) -- (1.250,-0.068) -- (1.235,-0.097) -- cycle;
\draw[ink2, line width=0.6pt, line cap=round] (0.592,-0.334) -- (0.604,-0.329);
\draw[ink2, line width=0.6pt, line cap=round] (1.235,-0.097) -- (1.250,-0.068);
\draw[loop, line cap=round] (1.074,-0.157) -- (1.089,-0.133);
\filldraw[fill=slab!36!paper, draw=slab!36!paper, line width=0.2pt] (-1.235,-0.485) -- (-1.219,-0.504) -- (-0.579,0.039) -- (-0.592,0.054) -- cycle;
\draw[ink2, line width=0.6pt, line cap=round] (-1.235,-0.485) -- (-1.219,-0.504);
\draw[ink2, line width=0.6pt, line cap=round] (-0.592,0.054) -- (-0.579,0.039);
\draw[loop, line cap=round] (-0.753,-0.081) -- (-0.739,-0.097);
\draw[loop, line cap=round] (-1.074,-0.350) -- (-1.059,-0.368);
\filldraw[fill=slab!54!paper, draw=slab!54!paper, line width=0.2pt] (0.579,-0.338) -- (0.592,-0.334) -- (1.235,-0.097) -- (1.219,-0.126) -- cycle;
\draw[ink2, line width=0.6pt, line cap=round] (0.579,-0.338) -- (0.592,-0.334);
\draw[ink2, line width=0.6pt, line cap=round] (1.219,-0.126) -- (1.235,-0.097);
\draw[loop, line cap=round] (1.059,-0.179) -- (1.074,-0.157);
\filldraw[fill=slab!36!paper, draw=slab!36!paper, line width=0.2pt] (-1.219,-0.504) -- (-1.201,-0.522) -- (-0.565,0.024) -- (-0.579,0.039) -- cycle;
\draw[ink2, line width=0.6pt, line cap=round] (-1.219,-0.504) -- (-1.201,-0.522);
\draw[ink2, line width=0.6pt, line cap=round] (-0.579,0.039) -- (-0.565,0.024);
\draw[loop, line cap=round] (-0.739,-0.097) -- (-0.724,-0.112);
\draw[loop, line cap=round] (-1.059,-0.368) -- (-1.042,-0.385);
\filldraw[fill=slab!54!paper, draw=slab!54!paper, line width=0.2pt] (0.565,-0.343) -- (0.579,-0.338) -- (1.219,-0.126) -- (1.201,-0.155) -- cycle;
\draw[ink2, line width=0.6pt, line cap=round] (0.565,-0.343) -- (0.579,-0.338);
\draw[ink2, line width=0.6pt, line cap=round] (1.201,-0.155) -- (1.219,-0.126);
\draw[loop, line cap=round] (1.042,-0.202) -- (1.059,-0.179);
\draw[loop, line cap=round] (0.724,-0.296) -- (0.739,-0.285);
\filldraw[fill=slab!35!paper, draw=slab!35!paper, line width=0.2pt] (-1.201,-0.522) -- (-1.181,-0.539) -- (-0.551,0.010) -- (-0.565,0.024) -- cycle;
\draw[ink2, line width=0.6pt, line cap=round] (-1.201,-0.522) -- (-1.181,-0.539);
\draw[ink2, line width=0.6pt, line cap=round] (-0.565,0.024) -- (-0.551,0.010);
\draw[loop, line cap=round] (-0.724,-0.112) -- (-0.709,-0.128);
\draw[loop, line cap=round] (-1.042,-0.385) -- (-1.024,-0.402);
\filldraw[fill=slab!54!paper, draw=slab!54!paper, line width=0.2pt] (0.551,-0.347) -- (0.565,-0.343) -- (1.201,-0.155) -- (1.181,-0.183) -- cycle;
\draw[ink2, line width=0.6pt, line cap=round] (0.551,-0.347) -- (0.565,-0.343);
\draw[ink2, line width=0.6pt, line cap=round] (1.181,-0.183) -- (1.201,-0.155);
\draw[ink3!50, line width=0.25pt] (0.551,-0.347) -- (1.181,-0.183);
\draw[loop, line cap=round] (1.024,-0.224) -- (1.042,-0.202);
\draw[loop, line cap=round] (0.709,-0.306) -- (0.724,-0.296);
\filldraw[fill=slab!35!paper, draw=slab!35!paper, line width=0.2pt] (-1.181,-0.539) -- (-1.160,-0.557) -- (-0.537,-0.005) -- (-0.551,0.010) -- cycle;
\draw[ink2, line width=0.6pt, line cap=round] (-1.181,-0.539) -- (-1.160,-0.557);
\draw[ink2, line width=0.6pt, line cap=round] (-0.551,0.010) -- (-0.537,-0.005);
\draw[ink3!50, line width=0.25pt] (-1.181,-0.539) -- (-0.551,0.010);
\draw[loop, line cap=round] (-0.709,-0.128) -- (-0.692,-0.143);
\draw[loop, line cap=round] (-1.024,-0.402) -- (-1.004,-0.419);
\filldraw[fill=slab!54!paper, draw=slab!54!paper, line width=0.2pt] (0.537,-0.350) -- (0.551,-0.347) -- (1.181,-0.183) -- (1.160,-0.211) -- cycle;
\draw[ink2, line width=0.6pt, line cap=round] (0.537,-0.350) -- (0.551,-0.347);
\draw[ink2, line width=0.6pt, line cap=round] (1.160,-0.211) -- (1.181,-0.183);
\draw[loop, line cap=round] (1.004,-0.246) -- (1.024,-0.224);
\draw[loop, line cap=round] (0.692,-0.316) -- (0.709,-0.306);
\filldraw[fill=slab!54!paper, draw=slab!54!paper, line width=0.2pt] (0.522,-0.354) -- (0.537,-0.350) -- (1.160,-0.211) -- (1.136,-0.239) -- cycle;
\draw[ink2, line width=0.6pt, line cap=round] (0.522,-0.354) -- (0.537,-0.350);
\draw[ink2, line width=0.6pt, line cap=round] (1.136,-0.239) -- (1.160,-0.211);
\draw[loop, line cap=round] (0.983,-0.268) -- (1.004,-0.246);
\draw[loop, line cap=round] (0.675,-0.325) -- (0.692,-0.316);
\filldraw[fill=slab!34!paper, draw=slab!34!paper, line width=0.2pt] (-1.160,-0.557) -- (-1.136,-0.573) -- (-0.522,-0.019) -- (-0.537,-0.005) -- cycle;
\draw[ink2, line width=0.6pt, line cap=round] (-1.160,-0.557) -- (-1.136,-0.573);
\draw[ink2, line width=0.6pt, line cap=round] (-0.537,-0.005) -- (-0.522,-0.019);
\draw[loop, line cap=round] (-0.692,-0.143) -- (-0.675,-0.158);
\filldraw[fill=slab!54!paper, draw=slab!54!paper, line width=0.2pt] (0.506,-0.357) -- (0.522,-0.354) -- (1.136,-0.239) -- (1.112,-0.266) -- cycle;
\draw[ink2, line width=0.6pt, line cap=round] (0.506,-0.357) -- (0.522,-0.354);
\draw[ink2, line width=0.6pt, line cap=round] (1.112,-0.266) -- (1.136,-0.239);
\draw[loop, line cap=round] (0.960,-0.289) -- (0.983,-0.268);
\filldraw[fill=slab!34!paper, draw=slab!34!paper, line width=0.2pt] (-1.136,-0.573) -- (-1.112,-0.590) -- (-0.506,-0.033) -- (-0.522,-0.019) -- cycle;
\draw[ink2, line width=0.6pt, line cap=round] (-1.136,-0.573) -- (-1.112,-0.590);
\draw[ink2, line width=0.6pt, line cap=round] (-0.522,-0.019) -- (-0.506,-0.033);
\draw[loop, line cap=round] (-0.675,-0.158) -- (-0.658,-0.172);
\filldraw[fill=slab!54!paper, draw=slab!54!paper, line width=0.2pt] (0.491,-0.359) -- (0.506,-0.357) -- (1.112,-0.266) -- (1.085,-0.293) -- cycle;
\draw[ink2, line width=0.6pt, line cap=round] (0.491,-0.359) -- (0.506,-0.357);
\draw[ink2, line width=0.6pt, line cap=round] (1.085,-0.293) -- (1.112,-0.266);
\draw[loop, line cap=round] (0.937,-0.310) -- (0.960,-0.289);
\filldraw[fill=slab!33!paper, draw=slab!33!paper, line width=0.2pt] (-1.112,-0.590) -- (-1.085,-0.606) -- (-0.491,-0.047) -- (-0.506,-0.033) -- cycle;
\draw[ink2, line width=0.6pt, line cap=round] (-1.112,-0.590) -- (-1.085,-0.606);
\draw[ink2, line width=0.6pt, line cap=round] (-0.506,-0.033) -- (-0.491,-0.047);
\draw[loop, line cap=round] (-0.658,-0.172) -- (-0.639,-0.187);
\filldraw[fill=slab!54!paper, draw=slab!54!paper, line width=0.2pt] (0.474,-0.361) -- (0.491,-0.359) -- (1.085,-0.293) -- (1.058,-0.320) -- cycle;
\draw[ink2, line width=0.6pt, line cap=round] (0.474,-0.361) -- (0.491,-0.359);
\draw[ink2, line width=0.6pt, line cap=round] (1.058,-0.320) -- (1.085,-0.293);
\draw[loop, line cap=round] (0.912,-0.330) -- (0.937,-0.310);
\filldraw[fill=slab!33!paper, draw=slab!33!paper, line width=0.2pt] (-1.085,-0.606) -- (-1.058,-0.621) -- (-0.474,-0.060) -- (-0.491,-0.047) -- cycle;
\draw[ink2, line width=0.6pt, line cap=round] (-1.085,-0.606) -- (-1.058,-0.621);
\draw[ink2, line width=0.6pt, line cap=round] (-0.491,-0.047) -- (-0.474,-0.060);
\draw[loop, line cap=round] (-0.639,-0.187) -- (-0.620,-0.201);
\draw[loop, line cap=round] (-0.937,-0.466) -- (-0.912,-0.481);
\filldraw[fill=slab!33!paper, draw=slab!33!paper, line width=0.2pt] (-1.058,-0.621) -- (-1.028,-0.636) -- (-0.458,-0.074) -- (-0.474,-0.060) -- cycle;
\draw[ink2, line width=0.6pt, line cap=round] (-1.058,-0.621) -- (-1.028,-0.636);
\draw[ink2, line width=0.6pt, line cap=round] (-0.474,-0.060) -- (-0.458,-0.074);
\draw[loop, line cap=round] (-0.620,-0.201) -- (-0.601,-0.214);
\draw[loop, line cap=round] (-0.912,-0.481) -- (-0.886,-0.495);
\filldraw[fill=slab!54!paper, draw=slab!54!paper, line width=0.2pt] (0.458,-0.363) -- (0.474,-0.361) -- (1.058,-0.320) -- (1.028,-0.346) -- cycle;
\draw[ink2, line width=0.6pt, line cap=round] (0.458,-0.363) -- (0.474,-0.361);
\draw[ink2, line width=0.6pt, line cap=round] (1.028,-0.346) -- (1.058,-0.320);
\draw[loop, line cap=round] (0.886,-0.350) -- (0.912,-0.330);
\draw[loop, line cap=round] (0.601,-0.359) -- (0.620,-0.351);
\filldraw[fill=slab!54!paper, draw=slab!54!paper, line width=0.2pt] (0.441,-0.365) -- (0.458,-0.363) -- (1.028,-0.346) -- (0.997,-0.371) -- cycle;
\filldraw[fill=slab!33!paper, draw=slab!33!paper, line width=0.2pt] (-1.028,-0.636) -- (-0.997,-0.650) -- (-0.441,-0.087) -- (-0.458,-0.074) -- cycle;
\draw[ink2, line width=0.6pt, line cap=round] (0.441,-0.365) -- (0.458,-0.363);
\draw[ink2, line width=0.6pt, line cap=round] (0.997,-0.371) -- (1.028,-0.346);
\draw[ink2, line width=0.6pt, line cap=round] (-1.028,-0.636) -- (-0.997,-0.650);
\draw[ink2, line width=0.6pt, line cap=round] (-0.458,-0.074) -- (-0.441,-0.087);
\draw[loop, line cap=round] (0.858,-0.370) -- (0.886,-0.350);
\draw[loop, line cap=round] (0.580,-0.367) -- (0.601,-0.359);
\draw[loop, line cap=round] (-0.601,-0.214) -- (-0.580,-0.227);
\draw[loop, line cap=round] (-0.886,-0.495) -- (-0.858,-0.509);
\filldraw[fill=slab!33!paper, draw=slab!33!paper, line width=0.2pt] (-0.997,-0.650) -- (-0.965,-0.663) -- (-0.424,-0.099) -- (-0.441,-0.087) -- cycle;
\draw[ink2, line width=0.6pt, line cap=round] (-0.997,-0.650) -- (-0.965,-0.663);
\draw[ink2, line width=0.6pt, line cap=round] (-0.441,-0.087) -- (-0.424,-0.099);
\draw[loop, line cap=round] (-0.580,-0.227) -- (-0.559,-0.240);
\draw[loop, line cap=round] (-0.858,-0.509) -- (-0.830,-0.522);
\filldraw[fill=slab!53!paper, draw=slab!53!paper, line width=0.2pt] (0.424,-0.366) -- (0.441,-0.365) -- (0.997,-0.371) -- (0.965,-0.396) -- cycle;
\draw[ink2, line width=0.6pt, line cap=round] (0.424,-0.366) -- (0.441,-0.365);
\draw[ink2, line width=0.6pt, line cap=round] (0.965,-0.396) -- (0.997,-0.371);
\draw[loop, line cap=round] (0.830,-0.389) -- (0.858,-0.370);
\draw[loop, line cap=round] (0.559,-0.374) -- (0.580,-0.367);
\filldraw[fill=slab!53!paper, draw=slab!53!paper, line width=0.2pt] (0.407,-0.367) -- (0.424,-0.366) -- (0.965,-0.396) -- (0.931,-0.421) -- cycle;
\filldraw[fill=slab!33!paper, draw=slab!33!paper, line width=0.2pt] (-0.965,-0.663) -- (-0.931,-0.676) -- (-0.407,-0.112) -- (-0.424,-0.099) -- cycle;
\draw[ink2, line width=0.6pt, line cap=round] (0.407,-0.367) -- (0.424,-0.366);
\draw[ink2, line width=0.6pt, line cap=round] (0.931,-0.421) -- (0.965,-0.396);
\draw[ink2, line width=0.6pt, line cap=round] (-0.965,-0.663) -- (-0.931,-0.676);
\draw[ink2, line width=0.6pt, line cap=round] (-0.424,-0.099) -- (-0.407,-0.112);
\draw[loop, line cap=round] (0.800,-0.407) -- (0.830,-0.389);
\draw[loop, line cap=round] (0.538,-0.380) -- (0.559,-0.374);
\draw[loop, line cap=round] (-0.559,-0.240) -- (-0.538,-0.253);
\filldraw[fill=slab!53!paper, draw=slab!53!paper, line width=0.2pt] (0.389,-0.368) -- (0.407,-0.367) -- (0.931,-0.421) -- (0.896,-0.444) -- cycle;
\filldraw[fill=slab!33!paper, draw=slab!33!paper, line width=0.2pt] (-0.931,-0.676) -- (-0.896,-0.688) -- (-0.389,-0.124) -- (-0.407,-0.112) -- cycle;
\draw[ink2, line width=0.6pt, line cap=round] (0.389,-0.368) -- (0.407,-0.367);
\draw[ink2, line width=0.6pt, line cap=round] (0.896,-0.444) -- (0.931,-0.421);
\draw[ink2, line width=0.6pt, line cap=round] (-0.931,-0.676) -- (-0.896,-0.688);
\draw[ink2, line width=0.6pt, line cap=round] (-0.407,-0.112) -- (-0.389,-0.124);
\draw[loop, line cap=round] (0.770,-0.425) -- (0.800,-0.407);
\draw[loop, line cap=round] (-0.538,-0.253) -- (-0.516,-0.265);
\filldraw[fill=slab!52!paper, draw=slab!52!paper, line width=0.2pt] (0.371,-0.368) -- (0.389,-0.368) -- (0.896,-0.444) -- (0.860,-0.467) -- cycle;
\filldraw[fill=slab!33!paper, draw=slab!33!paper, line width=0.2pt] (-0.896,-0.688) -- (-0.860,-0.699) -- (-0.371,-0.136) -- (-0.389,-0.124) -- cycle;
\draw[ink2, line width=0.6pt, line cap=round] (0.371,-0.368) -- (0.389,-0.368);
\draw[ink2, line width=0.6pt, line cap=round] (0.860,-0.467) -- (0.896,-0.444);
\draw[ink2, line width=0.6pt, line cap=round] (-0.896,-0.688) -- (-0.860,-0.699);
\draw[ink2, line width=0.6pt, line cap=round] (-0.389,-0.124) -- (-0.371,-0.136);
\draw[loop, line cap=round] (0.738,-0.443) -- (0.770,-0.425);
\draw[loop, line cap=round] (-0.516,-0.265) -- (-0.493,-0.277);
\filldraw[fill=slab!52!paper, draw=slab!52!paper, line width=0.2pt] (0.353,-0.368) -- (0.371,-0.368) -- (0.860,-0.467) -- (0.823,-0.490) -- cycle;
\filldraw[fill=slab!33!paper, draw=slab!33!paper, line width=0.2pt] (-0.860,-0.699) -- (-0.823,-0.710) -- (-0.353,-0.147) -- (-0.371,-0.136) -- cycle;
\draw[ink2, line width=0.6pt, line cap=round] (0.353,-0.368) -- (0.371,-0.368);
\draw[ink2, line width=0.6pt, line cap=round] (0.823,-0.490) -- (0.860,-0.467);
\draw[ink2, line width=0.6pt, line cap=round] (-0.860,-0.699) -- (-0.823,-0.710);
\draw[ink2, line width=0.6pt, line cap=round] (-0.371,-0.136) -- (-0.353,-0.147);
\draw[loop, line cap=round] (0.705,-0.459) -- (0.738,-0.443);
\draw[loop, line cap=round] (-0.493,-0.277) -- (-0.470,-0.288);
\draw[loop, line cap=round] (-0.738,-0.558) -- (-0.705,-0.569);
\filldraw[fill=slab!52!paper, draw=slab!52!paper, line width=0.2pt] (0.335,-0.367) -- (0.353,-0.368) -- (0.823,-0.490) -- (0.784,-0.512) -- cycle;
\draw[ink2, line width=0.6pt, line cap=round] (0.335,-0.367) -- (0.353,-0.368);
\draw[ink2, line width=0.6pt, line cap=round] (0.784,-0.512) -- (0.823,-0.490);
\draw[loop, line cap=round] (0.671,-0.475) -- (0.705,-0.459);
\draw[loop, line cap=round] (0.447,-0.403) -- (0.470,-0.398);
\filldraw[fill=slab!33!paper, draw=slab!33!paper, line width=0.2pt] (-0.823,-0.710) -- (-0.784,-0.720) -- (-0.335,-0.159) -- (-0.353,-0.147) -- cycle;
\draw[ink2, line width=0.6pt, line cap=round] (-0.823,-0.710) -- (-0.784,-0.720);
\draw[ink2, line width=0.6pt, line cap=round] (-0.353,-0.147) -- (-0.335,-0.159);
\draw[loop, line cap=round] (-0.470,-0.288) -- (-0.447,-0.299);
\draw[loop, line cap=round] (-0.705,-0.569) -- (-0.671,-0.580);
\filldraw[fill=slab!51!paper, draw=slab!51!paper, line width=0.2pt] (0.316,-0.366) -- (0.335,-0.367) -- (0.784,-0.512) -- (0.744,-0.533) -- cycle;
\draw[ink2, line width=0.6pt, line cap=round] (0.316,-0.366) -- (0.335,-0.367);
\draw[ink2, line width=0.6pt, line cap=round] (0.744,-0.533) -- (0.784,-0.512);
\draw[loop, line cap=round] (0.637,-0.491) -- (0.671,-0.475);
\draw[loop, line cap=round] (0.423,-0.408) -- (0.447,-0.403);
\filldraw[fill=slab!33!paper, draw=slab!33!paper, line width=0.2pt] (-0.784,-0.720) -- (-0.744,-0.729) -- (-0.316,-0.170) -- (-0.335,-0.159) -- cycle;
\draw[ink2, line width=0.6pt, line cap=round] (-0.784,-0.720) -- (-0.744,-0.729);
\draw[ink2, line width=0.6pt, line cap=round] (-0.335,-0.159) -- (-0.316,-0.170);
\draw[loop, line cap=round] (-0.447,-0.299) -- (-0.423,-0.310);
\draw[loop, line cap=round] (-0.671,-0.580) -- (-0.637,-0.589);
\filldraw[fill=slab!33!paper, draw=slab!33!paper, line width=0.2pt] (-0.744,-0.729) -- (-0.703,-0.737) -- (-0.297,-0.181) -- (-0.316,-0.170) -- cycle;
\draw[ink2, line width=0.6pt, line cap=round] (-0.744,-0.729) -- (-0.703,-0.737);
\draw[ink2, line width=0.6pt, line cap=round] (-0.316,-0.170) -- (-0.297,-0.181);
\draw[loop, line cap=round] (-0.423,-0.310) -- (-0.399,-0.320);
\draw[loop, line cap=round] (-0.637,-0.589) -- (-0.601,-0.598);
\filldraw[fill=slab!51!paper, draw=slab!51!paper, line width=0.2pt] (0.297,-0.365) -- (0.316,-0.366) -- (0.744,-0.533) -- (0.703,-0.553) -- cycle;
\draw[ink2, line width=0.6pt, line cap=round] (0.297,-0.365) -- (0.316,-0.366);
\draw[ink2, line width=0.6pt, line cap=round] (0.703,-0.553) -- (0.744,-0.533);
\draw[ink3!50, line width=0.25pt] (0.297,-0.365) -- (0.703,-0.553);
\draw[loop, line cap=round] (0.601,-0.506) -- (0.637,-0.491);
\draw[loop, line cap=round] (0.399,-0.412) -- (0.423,-0.408);
\filldraw[fill=slab!50!paper, draw=slab!50!paper, line width=0.2pt] (0.278,-0.363) -- (0.297,-0.365) -- (0.703,-0.553) -- (0.661,-0.572) -- cycle;
\filldraw[fill=slab!34!paper, draw=slab!34!paper, line width=0.2pt] (-0.703,-0.737) -- (-0.661,-0.745) -- (-0.278,-0.191) -- (-0.297,-0.181) -- cycle;
\draw[ink2, line width=0.6pt, line cap=round] (0.278,-0.363) -- (0.297,-0.365);
\draw[ink2, line width=0.6pt, line cap=round] (0.661,-0.572) -- (0.703,-0.553);
\draw[ink2, line width=0.6pt, line cap=round] (-0.703,-0.737) -- (-0.661,-0.745);
\draw[ink2, line width=0.6pt, line cap=round] (-0.297,-0.181) -- (-0.278,-0.191);
\draw[ink3!50, line width=0.25pt] (-0.703,-0.737) -- (-0.297,-0.181);
\draw[loop, line cap=round] (0.565,-0.520) -- (0.601,-0.506);
\draw[loop, line cap=round] (0.374,-0.416) -- (0.399,-0.412);
\draw[loop, line cap=round] (-0.399,-0.320) -- (-0.374,-0.329);
\filldraw[fill=slab!34!paper, draw=slab!34!paper, line width=0.2pt] (-0.661,-0.745) -- (-0.618,-0.751) -- (-0.259,-0.201) -- (-0.278,-0.191) -- cycle;
\draw[ink2, line width=0.6pt, line cap=round] (-0.661,-0.745) -- (-0.618,-0.751);
\draw[ink2, line width=0.6pt, line cap=round] (-0.278,-0.191) -- (-0.259,-0.201);
\draw[loop, line cap=round] (-0.374,-0.329) -- (-0.349,-0.339);
\filldraw[fill=slab!50!paper, draw=slab!50!paper, line width=0.2pt] (0.259,-0.361) -- (0.278,-0.363) -- (0.661,-0.572) -- (0.618,-0.591) -- cycle;
\draw[ink2, line width=0.6pt, line cap=round] (0.259,-0.361) -- (0.278,-0.363);
\draw[ink2, line width=0.6pt, line cap=round] (0.618,-0.591) -- (0.661,-0.572);
\draw[loop, line cap=round] (0.528,-0.534) -- (0.565,-0.520);
\filldraw[fill=slab!49!paper, draw=slab!49!paper, line width=0.2pt] (0.240,-0.359) -- (0.259,-0.361) -- (0.618,-0.591) -- (0.574,-0.609) -- cycle;
\filldraw[fill=slab!34!paper, draw=slab!34!paper, line width=0.2pt] (-0.618,-0.751) -- (-0.574,-0.757) -- (-0.240,-0.211) -- (-0.259,-0.201) -- cycle;
\draw[ink2, line width=0.6pt, line cap=round] (0.240,-0.359) -- (0.259,-0.361);
\draw[ink2, line width=0.6pt, line cap=round] (0.574,-0.609) -- (0.618,-0.591);
\draw[ink2, line width=0.6pt, line cap=round] (-0.618,-0.751) -- (-0.574,-0.757);
\draw[ink2, line width=0.6pt, line cap=round] (-0.259,-0.201) -- (-0.240,-0.211);
\draw[loop, line cap=round] (0.490,-0.547) -- (0.528,-0.534);
\draw[loop, line cap=round] (-0.349,-0.339) -- (-0.323,-0.348);
\filldraw[fill=slab!49!paper, draw=slab!49!paper, line width=0.2pt] (0.220,-0.357) -- (0.240,-0.359) -- (0.574,-0.609) -- (0.529,-0.626) -- cycle;
\draw[ink2, line width=0.6pt, line cap=round] (0.220,-0.357) -- (0.240,-0.359);
\draw[ink2, line width=0.6pt, line cap=round] (0.529,-0.626) -- (0.574,-0.609);
\draw[loop, line cap=round] (0.452,-0.559) -- (0.490,-0.547);
\filldraw[fill=slab!35!paper, draw=slab!35!paper, line width=0.2pt] (-0.574,-0.757) -- (-0.529,-0.762) -- (-0.220,-0.221) -- (-0.240,-0.211) -- cycle;
\draw[ink2, line width=0.6pt, line cap=round] (-0.574,-0.757) -- (-0.529,-0.762);
\draw[ink2, line width=0.6pt, line cap=round] (-0.240,-0.211) -- (-0.220,-0.221);
\draw[loop, line cap=round] (-0.323,-0.348) -- (-0.297,-0.356);
\draw[loop, line cap=round] (-0.490,-0.621) -- (-0.452,-0.627);
\filldraw[fill=slab!48!paper, draw=slab!48!paper, line width=0.2pt] (0.201,-0.354) -- (0.220,-0.357) -- (0.529,-0.626) -- (0.483,-0.642) -- cycle;
\draw[ink2, line width=0.6pt, line cap=round] (0.201,-0.354) -- (0.220,-0.357);
\draw[ink2, line width=0.6pt, line cap=round] (0.483,-0.642) -- (0.529,-0.626);
\draw[loop, line cap=round] (0.413,-0.570) -- (0.452,-0.559);
\draw[loop, line cap=round] (0.271,-0.426) -- (0.297,-0.424);
\filldraw[fill=slab!35!paper, draw=slab!35!paper, line width=0.2pt] (-0.529,-0.762) -- (-0.483,-0.766) -- (-0.201,-0.230) -- (-0.220,-0.221) -- cycle;
\draw[ink2, line width=0.6pt, line cap=round] (-0.529,-0.762) -- (-0.483,-0.766);
\draw[ink2, line width=0.6pt, line cap=round] (-0.220,-0.221) -- (-0.201,-0.230);
\draw[loop, line cap=round] (-0.297,-0.356) -- (-0.271,-0.364);
\draw[loop, line cap=round] (-0.452,-0.627) -- (-0.413,-0.632);
\filldraw[fill=slab!47!paper, draw=slab!47!paper, line width=0.2pt] (0.181,-0.351) -- (0.201,-0.354) -- (0.483,-0.642) -- (0.437,-0.657) -- cycle;
\filldraw[fill=slab!36!paper, draw=slab!36!paper, line width=0.2pt] (-0.483,-0.766) -- (-0.437,-0.769) -- (-0.181,-0.239) -- (-0.201,-0.230) -- cycle;
\draw[ink2, line width=0.6pt, line cap=round] (0.181,-0.351) -- (0.201,-0.354);
\draw[ink2, line width=0.6pt, line cap=round] (0.437,-0.657) -- (0.483,-0.642);
\draw[ink2, line width=0.6pt, line cap=round] (-0.483,-0.766) -- (-0.437,-0.769);
\draw[ink2, line width=0.6pt, line cap=round] (-0.201,-0.230) -- (-0.181,-0.239);
\draw[loop, line cap=round] (0.373,-0.581) -- (0.413,-0.570);
\draw[loop, line cap=round] (0.245,-0.427) -- (0.271,-0.426);
\draw[loop, line cap=round] (-0.271,-0.364) -- (-0.245,-0.372);
\draw[loop, line cap=round] (-0.413,-0.632) -- (-0.373,-0.636);
\filldraw[fill=slab!47!paper, draw=slab!47!paper, line width=0.2pt] (0.161,-0.347) -- (0.181,-0.351) -- (0.437,-0.657) -- (0.390,-0.672) -- cycle;
\draw[ink2, line width=0.6pt, line cap=round] (0.161,-0.347) -- (0.181,-0.351);
\draw[ink2, line width=0.6pt, line cap=round] (0.390,-0.672) -- (0.437,-0.657);
\draw[loop, line cap=round] (0.333,-0.591) -- (0.373,-0.581);
\draw[loop, line cap=round] (0.218,-0.428) -- (0.245,-0.427);
\filldraw[fill=slab!36!paper, draw=slab!36!paper, line width=0.2pt] (-0.437,-0.769) -- (-0.390,-0.771) -- (-0.161,-0.248) -- (-0.181,-0.239) -- cycle;
\draw[ink2, line width=0.6pt, line cap=round] (-0.437,-0.769) -- (-0.390,-0.771);
\draw[ink2, line width=0.6pt, line cap=round] (-0.181,-0.239) -- (-0.161,-0.248);
\draw[loop, line cap=round] (-0.245,-0.372) -- (-0.218,-0.379);
\draw[loop, line cap=round] (-0.373,-0.636) -- (-0.333,-0.640);
\filldraw[fill=slab!46!paper, draw=slab!46!paper, line width=0.2pt] (0.141,-0.343) -- (0.161,-0.347) -- (0.390,-0.672) -- (0.343,-0.685) -- cycle;
\draw[ink2, line width=0.6pt, line cap=round] (0.141,-0.343) -- (0.161,-0.347);
\draw[ink2, line width=0.6pt, line cap=round] (0.343,-0.685) -- (0.390,-0.672);
\draw[loop, line cap=round] (0.292,-0.600) -- (0.333,-0.591);
\draw[loop, line cap=round] (0.191,-0.429) -- (0.218,-0.428);
\filldraw[fill=slab!37!paper, draw=slab!37!paper, line width=0.2pt] (-0.390,-0.771) -- (-0.343,-0.772) -- (-0.141,-0.256) -- (-0.161,-0.248) -- cycle;
\draw[ink2, line width=0.6pt, line cap=round] (-0.390,-0.771) -- (-0.343,-0.772);
\draw[ink2, line width=0.6pt, line cap=round] (-0.161,-0.248) -- (-0.141,-0.256);
\draw[loop, line cap=round] (-0.218,-0.379) -- (-0.191,-0.385);
\filldraw[fill=slab!46!paper, draw=slab!46!paper, line width=0.2pt] (0.121,-0.339) -- (0.141,-0.343) -- (0.343,-0.685) -- (0.295,-0.698) -- cycle;
\filldraw[fill=slab!37!paper, draw=slab!37!paper, line width=0.2pt] (-0.343,-0.772) -- (-0.295,-0.772) -- (-0.121,-0.265) -- (-0.141,-0.256) -- cycle;
\draw[ink2, line width=0.6pt, line cap=round] (0.121,-0.339) -- (0.141,-0.343);
\draw[ink2, line width=0.6pt, line cap=round] (0.295,-0.698) -- (0.343,-0.685);
\draw[ink2, line width=0.6pt, line cap=round] (-0.343,-0.772) -- (-0.295,-0.772);
\draw[ink2, line width=0.6pt, line cap=round] (-0.141,-0.256) -- (-0.121,-0.265);
\draw[loop, line cap=round] (0.251,-0.608) -- (0.292,-0.600);
\draw[loop, line cap=round] (-0.191,-0.385) -- (-0.164,-0.391);
\filldraw[fill=slab!38!paper, draw=slab!38!paper, line width=0.2pt] (-0.295,-0.772) -- (-0.246,-0.771) -- (-0.101,-0.272) -- (-0.121,-0.265) -- cycle;
\draw[ink2, line width=0.6pt, line cap=round] (-0.295,-0.772) -- (-0.246,-0.771);
\draw[ink2, line width=0.6pt, line cap=round] (-0.121,-0.265) -- (-0.101,-0.272);
\draw[loop, line cap=round] (-0.164,-0.391) -- (-0.137,-0.397);
\filldraw[fill=slab!45!paper, draw=slab!45!paper, line width=0.2pt] (0.101,-0.335) -- (0.121,-0.339) -- (0.295,-0.698) -- (0.246,-0.709) -- cycle;
\draw[ink2, line width=0.6pt, line cap=round] (0.101,-0.335) -- (0.121,-0.339);
\draw[ink2, line width=0.6pt, line cap=round] (0.246,-0.709) -- (0.295,-0.698);
\draw[loop, line cap=round] (0.210,-0.616) -- (0.251,-0.608);
\filldraw[fill=slab!38!paper, draw=slab!38!paper, line width=0.2pt] (-0.246,-0.771) -- (-0.197,-0.769) -- (-0.081,-0.280) -- (-0.101,-0.272) -- cycle;
\draw[ink2, line width=0.6pt, line cap=round] (-0.246,-0.771) -- (-0.197,-0.769);
\draw[ink2, line width=0.6pt, line cap=round] (-0.101,-0.272) -- (-0.081,-0.280);
\draw[loop, line cap=round] (-0.137,-0.397) -- (-0.110,-0.402);
\draw[loop, line cap=round] (-0.210,-0.647) -- (-0.168,-0.647);
\filldraw[fill=slab!44!paper, draw=slab!44!paper, line width=0.2pt] (0.081,-0.330) -- (0.101,-0.335) -- (0.246,-0.709) -- (0.197,-0.720) -- cycle;
\draw[ink2, line width=0.6pt, line cap=round] (0.081,-0.330) -- (0.101,-0.335);
\draw[ink2, line width=0.6pt, line cap=round] (0.197,-0.720) -- (0.246,-0.709);
\draw[loop, line cap=round] (0.168,-0.622) -- (0.210,-0.616);
\filldraw[fill=slab!44!paper, draw=slab!44!paper, line width=0.2pt] (0.061,-0.325) -- (0.081,-0.330) -- (0.197,-0.720) -- (0.148,-0.729) -- cycle;
\filldraw[fill=slab!39!paper, draw=slab!39!paper, line width=0.2pt] (-0.197,-0.769) -- (-0.148,-0.767) -- (-0.061,-0.287) -- (-0.081,-0.280) -- cycle;
\draw[ink2, line width=0.6pt, line cap=round] (0.061,-0.325) -- (0.081,-0.330);
\draw[ink2, line width=0.6pt, line cap=round] (0.148,-0.729) -- (0.197,-0.720);
\draw[ink2, line width=0.6pt, line cap=round] (-0.197,-0.769) -- (-0.148,-0.767);
\draw[ink2, line width=0.6pt, line cap=round] (-0.081,-0.280) -- (-0.061,-0.287);
\draw[loop, line cap=round] (0.126,-0.628) -- (0.168,-0.622);
\draw[loop, line cap=round] (0.083,-0.426) -- (0.110,-0.427);
\draw[loop, line cap=round] (-0.110,-0.402) -- (-0.083,-0.407);
\draw[loop, line cap=round] (-0.168,-0.647) -- (-0.126,-0.647);
\filldraw[fill=slab!40!paper, draw=slab!40!paper, line width=0.2pt] (-0.148,-0.767) -- (-0.099,-0.763) -- (-0.040,-0.294) -- (-0.061,-0.287) -- cycle;
\draw[ink2, line width=0.6pt, line cap=round] (-0.148,-0.767) -- (-0.099,-0.763);
\draw[ink2, line width=0.6pt, line cap=round] (-0.061,-0.287) -- (-0.040,-0.294);
\draw[loop, line cap=round] (-0.083,-0.407) -- (-0.055,-0.411);
\draw[loop, line cap=round] (-0.126,-0.647) -- (-0.084,-0.646);
\filldraw[fill=slab!43!paper, draw=slab!43!paper, line width=0.2pt] (0.040,-0.319) -- (0.061,-0.325) -- (0.148,-0.729) -- (0.099,-0.738) -- cycle;
\draw[ink2, line width=0.6pt, line cap=round] (0.040,-0.319) -- (0.061,-0.325);
\draw[ink2, line width=0.6pt, line cap=round] (0.099,-0.738) -- (0.148,-0.729);
\draw[loop, line cap=round] (0.084,-0.633) -- (0.126,-0.628);
\draw[loop, line cap=round] (0.055,-0.424) -- (0.083,-0.426);
\filldraw[fill=slab!40!paper, draw=slab!40!paper, line width=0.2pt] (-0.099,-0.763) -- (-0.050,-0.758) -- (-0.020,-0.301) -- (-0.040,-0.294) -- cycle;
\draw[ink2, line width=0.6pt, line cap=round] (-0.099,-0.763) -- (-0.050,-0.758);
\draw[ink2, line width=0.6pt, line cap=round] (-0.040,-0.294) -- (-0.020,-0.301);
\draw[loop, line cap=round] (-0.055,-0.411) -- (-0.028,-0.415);
\draw[loop, line cap=round] (-0.084,-0.646) -- (-0.042,-0.644);
\filldraw[fill=slab!42!paper, draw=slab!42!paper, line width=0.2pt] (0.020,-0.313) -- (0.040,-0.319) -- (0.099,-0.738) -- (0.050,-0.746) -- cycle;
\draw[ink2, line width=0.6pt, line cap=round] (0.020,-0.313) -- (0.040,-0.319);
\draw[ink2, line width=0.6pt, line cap=round] (0.050,-0.746) -- (0.099,-0.738);
\draw[loop, line cap=round] (0.042,-0.638) -- (0.084,-0.633);
\draw[loop, line cap=round] (0.028,-0.422) -- (0.055,-0.424);
\filldraw[fill=slab!41!paper, draw=slab!41!paper, line width=0.2pt] (0.000,-0.307) -- (0.020,-0.313) -- (0.050,-0.746) -- (0.000,-0.752) -- cycle;
\filldraw[fill=slab!41!paper, draw=slab!41!paper, line width=0.2pt] (-0.050,-0.758) -- (-0.000,-0.752) -- (-0.000,-0.307) -- (-0.020,-0.301) -- cycle;
\draw[ink2, line width=0.6pt, line cap=round] (0.000,-0.307) -- (0.020,-0.313);
\draw[ink2, line width=0.6pt, line cap=round] (0.000,-0.752) -- (0.050,-0.746);
\draw[ink3!50, line width=0.25pt] (0.000,-0.307) -- (0.000,-0.752);
\draw[ink2, line width=0.6pt, line cap=round] (-0.050,-0.758) -- (-0.000,-0.752);
\draw[ink2, line width=0.6pt, line cap=round] (-0.020,-0.301) -- (-0.000,-0.307);
\draw[loop, line cap=round] (0.000,-0.641) -- (0.042,-0.638);
\draw[loop, line cap=round] (-0.000,-0.419) -- (0.028,-0.422);
\draw[loop, line cap=round] (-0.028,-0.415) -- (-0.000,-0.419);
\coordinate (mstart) at (0.000,-0.641);
\coordinate (mlap) at (-0.000,-0.419);
\coordinate (marrA) at (1.068,0.328);\coordinate (marrA2) at (1.046,0.358);
\coordinate (marrB) at (-1.068,0.040);\coordinate (marrB2) at (-1.089,0.015);

    \fill[ink] (mstart) circle (2.2pt);
    \filldraw[fill=paper, draw=ink, line width=0.8pt] (mlap) circle (2.0pt);
    \draw[ink, line width=1.5pt, ->] (marrA) -- (marrA2);
  \end{scope}
  \draw[loop, spin] (11.55,-2.15) circle (0.38);
  \fill[ink] ($(11.55,-2.15)+(-40:0.38)$) circle (1.6pt);
  \node[lab, anchor=north, align=center] at (11.55,-2.6) {plain loop:\\once round};
  \node[lab, anchor=north west, align=left, font=\tiny] at (10.4,-3.65) {● start\\○ after one trip:\\\ \ the other side\\solid: trip one\\dashed: trip two};
  \node[lab, anchor=north, align=center] at (8.7,-5.0) {two trips round\\to come home};

  % the sound
  \begin{axis}[snd, at={(12.85cm,-4.5cm)}, height=3.4cm, xmin=0, xmax=1000, ymax=5.6,
      xtick={120,360,600,840}, xlabel={frequency, Hz \quad (loop of T = 1/240 s)}]
    \foreach \k in {240,480,720} { \edef\t{\noexpand\draw[ink3!60, densely dashed, line width=0.4pt] (axis cs:\k,0) -- (axis cs:\k,5.0);} \t }
    \addplot[line width=1.0pt, ink] table[x=f, y=mag] {mlohs-comb.dat};
    \node[font=\tiny, text=ink2, anchor=south] at (axis cs:240,5.0) {240};
    \node[font=\tiny, text=ink2, anchor=south] at (axis cs:480,5.0) {480};
    \node[font=\tiny, text=ink2, anchor=south] at (axis cs:720,5.0) {720};
  \end{axis}
  \node[lab, anchor=north west, align=left] at (12.75,-5.25) {dashed: where a plain loop of T rings\\solid: the flipped loop rings between};
\end{scope}
\node[font=\scriptsize, anchor=north west, text=ink2, text width=19.6cm] at (0,-6.85) {A loop that flips its sign each trip must go round twice to return, so it rings at the odd harmonics of half its frequency. A clarinet's odd harmonics come the same way: the bore flips the wave's sign at its open end and not at the reed. Computed with a = 0.8, T = 200 samples at 48 kHz.};

% ---------- title block ----------
\draw[ink, line width=0.8pt] (-0.3,-7.95) rectangle (19.9,-8.85);
\draw[ink, line width=0.3pt] (9.6,-7.95) -- (9.6,-8.85) (14.6,-7.95) -- (14.6,-8.85);
\node[font=\scriptsize\bfseries, text=ink, anchor=west] at (-0.1,-8.4) {DSP IN LOOPING DIMENSIONS · THE MÖBIUS LOOP};
\node[font=\scriptsize, text=ink, anchor=west] at (9.8,-8.4) {INK computed};
\node[font=\scriptsize, text=ink2, anchor=west] at (14.8,-8.4) {FIG · 2026-11-09 · PAPER};

\begin{scope}[on background layer]
  \fill[paper] ($(current bounding box.south west)+(-0.4,-0.4)$) rectangle ($(current bounding box.north east)+(0.4,0.4)$);
\end{scope}
\end{tikzpicture}
\end{document}
